\documentclass[11pt]{article}

\usepackage[margin=1in]{geometry}
\usepackage{amsmath,amssymb,amsthm,mathtools}
\usepackage{enumitem}
\usepackage{booktabs}
\usepackage{array}
\usepackage{xcolor}
\usepackage{hyperref}
\usepackage{longtable}
\usepackage{tikz}

\hypersetup{colorlinks=true, linkcolor=blue, citecolor=blue, urlcolor=blue}

% ---------- Theorem environments ----------
\theoremstyle{definition}
\newtheorem{definition}{Definition}[section]
\newtheorem{example}[definition]{Example}
\newtheorem{problem}[definition]{Problem}

\theoremstyle{plain}
\newtheorem{theorem}[definition]{Theorem}
\newtheorem{proposition}[definition]{Proposition}
\newtheorem{lemma}[definition]{Lemma}
\newtheorem{corollary}[definition]{Corollary}
\newtheorem{axiom}[definition]{Axiom}

\theoremstyle{remark}
\newtheorem{remark}[definition]{Remark}
\newtheorem{intuition}[definition]{Intuition}

% ---------- Supplementary / "further reading" material box (RED) ----------
\newenvironment{extra}[1][Supplementary Material (beyond the core syllabus)]{%
  \begin{quote}\color{red}\noindent\textbf{#1.}\par\smallskip
}{%
  \end{quote}
}

% ---------- Custom macros ----------
\newcommand{\Prob}[1]{\mathbf{P}\!\left[#1\right]}
\newcommand{\E}[1]{\mathbf{E}\!\left[#1\right]}
\newcommand{\Var}[1]{\operatorname{Var}\!\left[#1\right]}
\newcommand{\SD}[1]{\operatorname{SD}\!\left[#1\right]}
\newcommand{\R}{\mathbb{R}}
\newcommand{\N}{\mathbb{N}}
\newcommand{\Nz}{\mathbb{N}_0}
\newcommand{\Z}{\mathbb{Z}}
\newcommand{\Om}{\Omega}
\newcommand{\iid}{\text{i.i.d.}}
\newcommand{\perm}[2]{{}^{#1}P_{#2}}
\newcommand{\multinom}[2]{\dbinom{#1}{#2}}
\newcommand{\dd}{\,d}

\title{\textbf{Rigorous Study Notes: Combinatorics and Probability Theory}\\
\large Compiled and expanded from lecture notes (5 August -- 10 September 2026)}
\author{}
\date{}

\begin{document}
\maketitle
\tableofcontents
\newpage

% =========================================================================
\section{Notation and Terminology: A Key}
% =========================================================================

Before proceeding, we fix the notation and vocabulary used throughout. Every symbol below is used consistently in the sections that follow.

\subsection{Set-theoretic and logical symbols}

\begin{center}
\renewcommand{\arraystretch}{1.3}
\begin{longtable}{p{2.7cm} p{11.5cm}}
\toprule
\textbf{Symbol} & \textbf{Meaning} \\
\midrule
$\in$, $\notin$ & ``is an element of'', ``is not an element of'' \\
$\subset$, $\subseteq$ & ``is a subset of'' (we use $\subset$ to allow equality unless stated otherwise, matching the source notes) \\
$\cup$, $\cap$ & union, intersection of sets \\
$A \setminus B$ & set difference: elements of $A$ not in $B$, i.e.\ $A \cap B^c$ \\
$A^c$ & complement of $A$ relative to the ambient sample space $S$ \\
$\varnothing$ & the empty set \\
$|A|$ & cardinality (number of elements) of a finite set $A$ \\
$2^{\Omega}$ & the power set of $\Omega$: the set of \emph{all} subsets of $\Omega$ \\
$\mathbb{N}$, $\mathbb{N}_0$ & $\{1,2,3,\dots\}$ and $\{0,1,2,3,\dots\}$ respectively \\
$\mathbb{R}$ & the real numbers \\
$\forall$, $\exists$ & ``for all'', ``there exists'' \\
s.t. & ``such that'' \\
$\iff$ & ``if and only if'' \\
$:=$ & ``is defined to be'' \\
$\{x_n\}_{n\in\N}$ & a sequence indexed by $n$ \\
$x_n \to x$ & the sequence $(x_n)$ converges to $x$ \\
$x_n \downarrow a$ & $(x_n)$ decreases monotonically to $a$ \\
\bottomrule
\end{longtable}
\end{center}

\subsection{Combinatorial and probabilistic symbols}

\begin{center}
\renewcommand{\arraystretch}{1.35}
\begin{longtable}{p{3.2cm} p{11cm}}
\toprule
\textbf{Symbol} & \textbf{Meaning} \\
\midrule
$n!$ & $n$ factorial, $n(n-1)(n-2)\cdots 1$, with $0! := 1$ \\
$\perm{n}{r}$ & number of \emph{permutations} of $r$ objects chosen from $n$ distinguishable objects, $= \dfrac{n!}{(n-r)!}$ \\
$\dbinom{n}{r}$ & the binomial coefficient, ``$n$ choose $r$'', $=\dfrac{n!}{r!(n-r)!}$: the number of $r$-element subsets of an $n$-element set \\
$\dbinom{n}{m_1,m_2,\dots,m_r}$ & the multinomial coefficient $=\dfrac{n!}{m_1!m_2!\cdots m_r!}$, where $m_1+\cdots+m_r=n$ \\
$S$ or $\Omega$ & the sample space of a random experiment: the set of all possible outcomes \\
$\omega$ & a generic sample point, $\omega \in S$ \\
$E, F, G, A, B, \dots$ & events, i.e.\ subsets of the sample space \\
$\mathbf{P}[E]$ or $\mathbf{P}(E)$ & the probability of the event $E$ \\
$\mathbf{P}[E \mid F]$ & the conditional probability of $E$ given $F$ \\
$X, Y, Z$ & random variables \\
$p(x)$ or $p_X(x)$ & the probability mass function (pmf) of a discrete random variable $X$: $p(x) = \mathbf{P}[X=x]$ \\
$f(x)$ or $f_X(x)$ & the probability density function (pdf) of a continuous random variable $X$ \\
$F(x)$ or $F_X(x)$ & the cumulative distribution function (cdf): $F(x) = \mathbf{P}[X \le x]$ \\
$\mathbf{E}[X]$ & the expectation (mean) of $X$ \\
$\operatorname{Var}[X]$ & the variance of $X$ \\
$\operatorname{SD}[X]$ & the standard deviation of $X$, $=\sqrt{\operatorname{Var}[X]}$ \\
$\mathbf{1}_A$ & the indicator random variable of an event $A$: equals $1$ if $A$ occurs, $0$ otherwise \\
$X \sim \text{Binomial}(n,p)$ & ``$X$ is distributed as a Binomial random variable with parameters $n,p$'' \\
i.i.d. & ``independent and identically distributed'' \\
$q$ & shorthand for $1-p$, when $p$ denotes a success probability \\
\bottomrule
\end{longtable}
\end{center}

\subsection{Key terminology (informal glossary)}

\begin{itemize}[leftmargin=1.6em]
\item \textbf{Random experiment}: a procedure whose set of possible outcomes is known in advance, but whose actual outcome, on any particular performance, cannot be predicted with certainty.
\item \textbf{Sample space}: the set of \emph{all} possible outcomes of a random experiment.
\item \textbf{Sample point}: a single element of the sample space, i.e.\ one possible outcome.
\item \textbf{Event}: any subset of the sample space; informally, a statement about the outcome that is either true or false once the experiment is performed.
\item \textbf{Mutually exclusive (disjoint) events}: events $E,F$ with $E \cap F = \varnothing$; they cannot both occur.
\item \textbf{Random variable}: a real-valued (or vector-valued) function defined on the sample space.
\item \textbf{Discrete random variable}: one taking at most countably many values.
\item \textbf{Continuous random variable}: one whose cdf is (absolutely) continuous everywhere; equivalently, one possessing a probability density function.
\item \textbf{Support}: the set of values at which the pmf (or pdf) is (or can be) strictly positive.
\item \textbf{Independence}: informally, the property that knowledge of one event/random variable provides no information about another.
\end{itemize}

\bigskip
\textit{A note on rigor and intuition.} Every claim below that can be proved is proved --
using only the axioms and previously established results, in the style of the original notes.
Where a passage is genuinely subtle (for instance, the definition of a probability density, or
why exchangeability holds), we add a boxed \textbf{Intuition} paragraph with a concrete example.
Passages in \textcolor{red}{red} correspond to ``further reading'' material that goes strictly beyond
the examinable syllabus (typically requiring ideas from measure theory), included here in full
rigor for completeness.

\newpage
% =========================================================================
\section{Combinatorial Foundations}
% =========================================================================

Combinatorics is logically prior to probability theory because the \emph{naive} definition of
probability (Section~\ref{sec:naive-prob}) reduces every probability computation, on a finite
sample space with equally likely outcomes, to a pair of counting problems: counting the
favourable outcomes, and counting the total outcomes.

\subsection{The basic principle of counting}

\begin{proposition}[Basic principle of counting -- two experiments]
\label{prop:basic-counting-2}
If experiment $1$ can result in any one of $m$ possible outcomes, and, regardless of the outcome
of experiment $1$, experiment $2$ can result in any one of $n$ possible outcomes, then the two
experiments together can result in $mn$ possible outcome-pairs.
\end{proposition}

\begin{proof}
List the outcomes of experiment $1$ as $O_1, O_2, \dots, O_m$. For each fixed $O_i$, the possible
outcome-pairs with first coordinate $O_i$ are $(O_i, 1), (O_i,2), \dots, (O_i, n)$ -- exactly $n$
of them, one for each outcome of experiment $2$. These $m$ collections of pairs (one collection
per choice of $i$) are pairwise disjoint (a pair with first coordinate $O_i$ cannot equal a pair
with first coordinate $O_j$ for $i \ne j$), and their union is the set of \emph{all}
outcome-pairs. Hence the total number of pairs is
$\underbrace{n + n + \cdots + n}_{m \text{ times}} = mn$.
\end{proof}

\begin{proposition}[Generalized basic principle of counting]
If $r$ experiments are performed, and the $i$-th experiment (for $1 \le i \le r$) can result in
$m_i$ possible outcomes irrespective of the outcomes of the other experiments, then the total
number of possible $r$-tuples of outcomes is $m_1 m_2 \cdots m_r$.
\end{proposition}

\begin{proof}
Induction on $r$. The base case $r=1$ is trivial. Suppose the claim holds for $r-1$ experiments,
so that there are $m_1 m_2 \cdots m_{r-1}$ possible $(r-1)$-tuples from the first $r-1$
experiments. Now regard the ``first $r-1$ experiments taken together'' as a single experiment
with $m_1 \cdots m_{r-1}$ possible outcomes, and the $r$-th experiment as a second experiment
with $m_r$ outcomes. By Proposition~\ref{prop:basic-counting-2}, the total number of pairs
(($r-1$)-tuple, outcome of experiment $r$) -- which is precisely the number of $r$-tuples -- is
$(m_1 \cdots m_{r-1}) \cdot m_r = m_1 m_2 \cdots m_r$.
\end{proof}

\subsection{Permutations}

\begin{theorem}[Permutations of distinguishable objects]
The number of ways to arrange $n$ distinguishable objects in a row (i.e.\ the number of
permutations of $n$ distinct objects) is $n!$.
\end{theorem}

\begin{proof}
Think of filling $n$ positions, one at a time, from left to right. There are $n$ choices for
which object occupies position $1$. Having filled position $1$, there are $n-1$ remaining objects,
hence $n-1$ choices for position $2$; and so on, until there is exactly $1$ choice remaining for
position $n$. By the generalized basic principle of counting, the total number of arrangements is
$n(n-1)(n-2)\cdots 1 = n!$.
\end{proof}

\begin{theorem}[Permutations with repetition]
\label{thm:perm-repetition}
Suppose $n$ objects consist of $r$ distinguishable \emph{types}, with $k_1$ objects of type $1$,
$k_2$ of type $2$, \dots, $k_r$ of type $r$ (objects of the same type being mutually
indistinguishable), where $k_1 + k_2 + \cdots + k_r = n$. Then the number of distinguishable
arrangements of these $n$ objects is
\[
\frac{n!}{k_1!\,k_2! \cdots k_r!}.
\]
\end{theorem}

\begin{proof}
Temporarily label the $k_i$ objects of type $i$ as $i_1, i_2, \dots, i_{k_i}$, rendering all $n$
objects distinguishable. There are now $n!$ arrangements of these fully-labeled objects. Every
\emph{genuine} (unlabeled) arrangement of the original $n$ objects corresponds to exactly
$k_1! k_2! \cdots k_r!$ labeled arrangements, because within each block of $k_i$ same-type objects
the labels can be permuted amongst themselves in $k_i!$ ways without altering the underlying
unlabeled arrangement, and this can be done independently for each of the $r$ types (basic
principle of counting again). Hence the labeled arrangements are partitioned into classes of size
$k_1!\cdots k_r!$, each class corresponding to one genuine arrangement, so the number of genuine
arrangements is $n! / (k_1! \cdots k_r!)$.
\end{proof}

\begin{intuition}
Take $n=3$ letters with two \texttt{A}'s and one \texttt{B} ($k_1=2, k_2=1$). Labeling gives
$A_1 A_2 B, A_2 A_1 B, A_1 B A_2, \dots$ -- $3! = 6$ labeled strings. But $A_1A_2B$ and $A_2A_1B$
are the \emph{same} unlabeled string \texttt{AAB}. Indeed the genuine arrangements are
$\texttt{AAB}, \texttt{ABA}, \texttt{BAA}$ -- exactly $3!/(2!1!) = 3$ of them.
\end{intuition}

\subsection{Combinations}

\begin{definition}
Given $n$ distinguishable objects and $0 \le r \le n$, a \textbf{combination} of size $r$ is an
(unordered) subset of $r$ objects chosen from the $n$. We write $\dbinom{n}{r}$ for the number of
such subsets.
\end{definition}

\begin{theorem}
\label{thm:binom-coeff}
\[
\binom{n}{r} = \frac{n!}{(n-r)!\,r!} = \frac{n(n-1)\cdots(n-r+1)}{r!}.
\]
\end{theorem}

\begin{proof}[Proof via equivalence classes]
Consider the set $\mathcal{P}_r$ of all \emph{ordered} selections of $r$ distinct objects from the
$n$ objects -- these are the $r$-permutations, and by the same slot-filling argument as above,
$|\mathcal{P}_r| = n(n-1)\cdots(n-r+1) = \dfrac{n!}{(n-r)!}$.

Define a relation $\sim$ on $\mathcal{P}_r$: two ordered $r$-tuples of distinct objects are related
if they consist of exactly the same underlying \emph{set} of $r$ objects (differing only, perhaps,
in the order in which the objects are listed). One checks immediately that $\sim$ is reflexive,
symmetric, and transitive, i.e.\ an equivalence relation; it therefore partitions $\mathcal{P}_r$
into disjoint equivalence classes.

Fix any equivalence class: it consists precisely of all the orderings of one particular
$r$-element subset, and there are $r!$ such orderings (permutations of $r$ distinct objects).
Hence every equivalence class has exactly $r!$ elements. Moreover, the equivalence classes are in
bijection with the $r$-element subsets of the original $n$ objects (each class $\leftrightarrow$
the unique subset underlying it).

Since $\mathcal{P}_r$ is partitioned into equally-sized (size $r!$) classes, and the number of
classes equals $\dbinom{n}{r}$ (the quantity we want), we obtain
\[
|\mathcal{P}_r| = \binom{n}{r} \cdot r! \quad\Longrightarrow\quad \binom{n}{r} =
\frac{|\mathcal{P}_r|}{r!} = \frac{n!}{(n-r)!\,r!}. \qedhere
\]
\end{proof}

\begin{theorem}[Pascal's identity]
For $1 \le r \le n-1$,
\[
\binom{n}{r} = \binom{n-1}{r} + \binom{n-1}{r-1}.
\]
\end{theorem}

\begin{proof}[Combinatorial proof]
We count the $r$-element subsets of an $n$-element set $\{1,2,\dots,n\}$ in two ways. Fix the
element $n$. Every $r$-subset either
\begin{itemize}
\item does \emph{not} contain $n$: such subsets are exactly the $r$-subsets of $\{1,\dots,n-1\}$,
of which there are $\dbinom{n-1}{r}$; or
\item \emph{does} contain $n$: removing $n$ leaves an $(r-1)$-subset of $\{1,\dots,n-1\}$, and
this correspondence is a bijection, so there are $\dbinom{n-1}{r-1}$ such subsets.
\end{itemize}
These two cases are exhaustive and mutually exclusive, so by the addition of disjoint counts,
$\dbinom{n}{r} = \dbinom{n-1}{r} + \dbinom{n-1}{r-1}$.
\end{proof}

\begin{theorem}[Binomial theorem]
\label{thm:binomial-thm}
For any $n \in \Nz$,
\[
(x+y)^n = \sum_{k=0}^n \binom{n}{k} x^k y^{n-k}.
\]
\end{theorem}

\begin{proof}[Proof by induction on $n$]
\textbf{Base case} ($n=0$): both sides equal $1$.

\textbf{Inductive step}: suppose the identity holds for $n$. Then
\[
(x+y)^{n+1} = (x+y)(x+y)^n = (x+y)\sum_{k=0}^n \binom{n}{k}x^k y^{n-k}
= \sum_{k=0}^n \binom{n}{k}x^{k+1}y^{n-k} + \sum_{k=0}^n \binom{n}{k}x^k y^{n-k+1}.
\]
In the first sum, substitute $j = k+1$ (so $j$ ranges over $1,\dots,n+1$):
\[
\sum_{j=1}^{n+1}\binom{n}{j-1}x^j y^{n+1-j}.
\]
Relabel $j$ as $k$ in both sums and adopt the convention $\dbinom{n}{-1}=\dbinom{n}{n+1}=0$, so
both sums may now be taken over $k=0,\dots,n+1$:
\[
(x+y)^{n+1} = \sum_{k=0}^{n+1}\left[\binom{n}{k-1}+\binom{n}{k}\right]x^k y^{n+1-k}
= \sum_{k=0}^{n+1}\binom{n+1}{k}x^k y^{n+1-k},
\]
where the last equality is Pascal's identity. This completes the induction.
\end{proof}

\begin{proof}[Combinatorial proof]
Write $(x+y)^n$ as the product of $n$ factors, $(x+y)(x+y)\cdots(x+y)$. Expanding, each term of
the expansion is obtained by choosing, from \emph{each} of the $n$ factors, either the term $x$
or the term $y$, and multiplying the $n$ choices together. A term equal to $x^k y^{n-k}$ arises
precisely from those choices in which $x$ was selected from exactly $k$ of the $n$ factors (and
$y$ from the remaining $n-k$). The number of ways to choose \emph{which} $k$ of the $n$ factors
contribute an $x$ is $\dbinom{n}{k}$ (Theorem~\ref{thm:binom-coeff}). Summing over all
$k=0,\dots,n$ gives the stated identity.
\end{proof}

\subsection{Multinomial coefficients and the multinomial theorem}

\begin{theorem}[Multinomial coefficients]
Suppose $n$ \emph{distinct} objects are to be divided into $r$ \emph{distinguishable} (labeled)
groups, the $i$-th group required to contain exactly $m_i$ objects, where $m_1+\cdots+m_r=n$. The
number of ways to do this is
\[
\binom{n}{m_1,m_2,\dots,m_r} := \frac{n!}{m_1!\,m_2!\cdots m_r!}.
\]
\end{theorem}

\begin{proof}
Build the partition group by group. Choose the $m_1$ objects for group $1$ out of $n$:
$\dbinom{n}{m_1}$ ways. Then choose the $m_2$ objects for group $2$ out of the remaining $n-m_1$:
$\dbinom{n-m_1}{m_2}$ ways. Continue; at the $i$-th step there are $n - m_1 - \cdots - m_{i-1}$
objects left, and we choose $\dbinom{n-m_1-\cdots-m_{i-1}}{m_i}$ ways to select group $i$. By the
basic principle of counting, the total number of ways is the product
\[
\binom{n}{m_1}\binom{n-m_1}{m_2}\binom{n-m_1-m_2}{m_3}\cdots\binom{m_r}{m_r}.
\]
Writing out factorials, this telescopes:
\[
\frac{n!}{m_1!(n-m_1)!}\cdot\frac{(n-m_1)!}{m_2!(n-m_1-m_2)!}\cdot
\frac{(n-m_1-m_2)!}{m_3!(n-m_1-m_2-m_3)!}\cdots = \frac{n!}{m_1!m_2!\cdots m_r!}. \qedhere
\]
\end{proof}

\begin{remark}
Compare this with Theorem~\ref{thm:perm-repetition}. The two formulas coincide algebraically, but
the combinatorial \emph{interpretations} differ: Theorem~\ref{thm:perm-repetition} counts
arrangements of $n$ objects of $r$ indistinguishable types in a row, while the multinomial
coefficient counts assignments of $n$ distinct objects to $r$ labeled groups of prescribed sizes.
There is a natural bijection between the two: reading off, in a permutation-with-repetition
arrangement, \emph{which positions} hold type-$i$ objects, for each $i$, exactly recovers a
partition of the $n$ positions (now playing the role of ``$n$ distinct objects'') into $r$ labeled
groups of sizes $k_1,\dots,k_r$.
\end{remark}

\begin{theorem}[Multinomial theorem]
For any $n \in \Nz$ and $r \ge 1$,
\[
(x_1+x_2+\cdots+x_r)^n = \sum_{\substack{m_1,\dots,m_r \in \Nz \\ m_1+\cdots+m_r = n}}
\binom{n}{m_1,\dots,m_r} x_1^{m_1}x_2^{m_2}\cdots x_r^{m_r}.
\]
\end{theorem}

\begin{proof}
Exactly as in the combinatorial proof of the binomial theorem: write the left side as a product
of $n$ factors $(x_1+\cdots+x_r)$. A term $x_1^{m_1}\cdots x_r^{m_r}$ (with $\sum m_i = n$) arises
from those selections in which, for each $i$, the value $x_i$ was chosen from exactly $m_i$ of the
$n$ factors. Choosing, for each $i=1,\dots,r$, which $m_i$ of the $n$ factors contribute $x_i$ is
precisely an assignment of the $n$ (distinct, because we track \emph{which} factor) factors into
$r$ labeled groups of sizes $m_1,\dots,m_r$ -- there are $\dbinom{n}{m_1,\dots,m_r}$ such
assignments. Summing over all admissible $(m_1,\dots,m_r)$ gives the result.
\end{proof}

\subsection{Distributing balls into boxes (``stars and bars'')}

\begin{theorem}
\label{thm:stars-bars}
The number of ways to place $n$ indistinguishable balls into $k$ distinguishable (labeled) boxes
is
\[
\binom{n+k-1}{k-1} = \binom{n+k-1}{n}.
\]
If, in addition, each box is required to contain \emph{at least one} ball (so necessarily
$n \ge k$), the number of ways is
\[
\binom{n-1}{k-1}.
\]
\end{theorem}

\begin{proof}
\textbf{Unrestricted case.} Represent a placement of $n$ identical balls into $k$ labeled boxes
by a string of $n$ stars ($\star$) and $k-1$ bars ($|$): the stars before the first bar denote the
number of balls in box $1$, the stars between the first and second bars denote the number of
balls in box $2$, and so on, with the stars after the last bar denoting the number of balls in box
$k$. Every such string of $n$ stars and $k-1$ bars (in any order) corresponds to exactly one
placement, and conversely. The number of such strings is the number of ways to choose the
positions of the $k-1$ bars among the $n+k-1$ total symbols, i.e.\ $\dbinom{n+k-1}{k-1}$.

\textbf{At-least-one-ball case.} First place one ball in each of the $k$ boxes (using up $k$ of
the $n$ balls, forced and with no remaining choice). The remaining $n-k$ balls are then
distributed \emph{without restriction} among the $k$ boxes, which by the unrestricted case has
$\dbinom{(n-k)+k-1}{k-1} = \dbinom{n-1}{k-1}$ solutions.
\end{proof}

\begin{intuition}
This is exactly the combinatorics of writing a positive integer $n$ as an \emph{ordered} sum of
$k$ non-negative (respectively, positive) integers: $n = x_1 + x_2 + \cdots + x_k$ with $x_i \ge 0$
(resp. $x_i \ge 1$), where $x_i$ is the number of balls in box $i$.
\end{intuition}

\subsection{Worked example: antennas with no two defectives adjacent}

\begin{problem}
There are $n$ antennas in total, of which $m$ are defective, $m \le n$. All defective antennas are
mutually indistinguishable, and all functioning antennas are mutually indistinguishable. In how
many ways can the antennas be arranged in a row so that no two defective antennas are placed
consecutively?
\end{problem}

\begin{proof}[Solution]
First arrange the $n-m$ functioning antennas in a row; since they are indistinguishable, there is
exactly $\boxed{1}$ way to do this, and it creates $n-m+1$ ``gaps'' -- one before the first
functioning antenna, one after each functioning antenna, and one after the last:
\[
\_\, F \,\_\, F \,\_\, F \,\cdots\, F \,\_
\]
There are $n-m+1$ underscores (gaps). To ensure no two defective antennas are consecutive, each
gap may receive \emph{at most one} defective antenna (placing two or more defectives in the same
gap would put them adjacent to each other). We must choose $m$ of these $n-m+1$ gaps to each
receive exactly one (indistinguishable) defective antenna. The number of ways to choose which $m$
gaps receive a defective antenna is
\[
\binom{n-m+1}{m},
\]
which is therefore the number of valid arrangements (this requires $m \le n-m+1$, i.e.\
$m \le \lceil (n+1)/2 \rceil$, else the answer is $0$).
\end{proof}

\subsection{Summary comparison table}

\begin{center}
\renewcommand{\arraystretch}{1.5}
\begin{longtable}{p{4.3cm} p{3.7cm} p{2.2cm} p{4.3cm}}
\toprule
\textbf{Scenario} & \textbf{Formula} & \textbf{Order matters?} & \textbf{Repetition?} \\
\midrule
Arrangements of $n$ distinct objects & $n!$ & Yes & No \\
$r$-permutations of $n$ distinct objects & $\dfrac{n!}{(n-r)!}$ & Yes & No \\
$r$-combinations of $n$ distinct objects & $\dbinom{n}{r}$ & No & No \\
Arrangements of $n$ objects, $r$ types with multiplicities $k_1,\dots,k_r$ & $\dfrac{n!}{k_1!\cdots k_r!}$ & Yes (positions) & Yes (within type) \\
Partition $n$ distinct objects into $r$ labeled groups of sizes $m_1,\dots,m_r$ & $\dbinom{n}{m_1,\dots,m_r}$ & No (within group) & --- \\
$n$ identical balls into $k$ labeled boxes & $\dbinom{n+k-1}{k-1}$ & --- & Yes \\
$n$ identical balls into $k$ labeled boxes, each $\ge 1$ & $\dbinom{n-1}{k-1}$ & --- & Yes \\
\bottomrule
\end{longtable}
\end{center}

\newpage
% =========================================================================
\section{Foundations of Probability Theory}
% =========================================================================

\subsection{Random experiments, sample spaces, and events}

\begin{definition}
A \textbf{random experiment} is a procedure for which the set of all possible outcomes is known
in advance, but the outcome that will actually be observed on any given performance cannot be
predicted with certainty.
\end{definition}

\begin{example}
Administering a blood-pressure-lowering drug to a patient and measuring the resulting blood
pressure (in whole mmHg) is a random experiment: the set of conceivable readings is known (some
range of positive integers), but the actual reading cannot be known in advance, since it depends
on age, dosage, co-morbidities, allergies, and so on.
\end{example}

\begin{definition}
The \textbf{sample space} $S$ (also written $\Omega$) of a random experiment is the set of all
its possible outcomes. Each element $\omega \in S$ is a \textbf{sample point}.
\end{definition}

\begin{definition}
An \textbf{event} is any subset $E \subset S$ of the sample space.
\end{definition}

\begin{example}
Tossing a fair coin $3$ times gives sample space
\[
S = \{HHH, HHT, HTH, THH, HTT, THT, TTH, TTT\}.
\]
Throwing a die twice gives $S = \{(i,j) : 1 \le i,j \le 6\}$, $|S|=36$.
The event ``the sum of the two throws is even'' is
$E = \{(i,j) \in S : i+j \text{ even}\}$; the event ``the sum is divisible by $3$'' is
$F = \{(i,j)\in S : 3 \mid (i+j)\}$.
\end{example}

\subsection{The algebra of events}

Since events are sets, they obey the ordinary laws of set algebra. Given events $E,F,G$ (subsets
of the same $S$):
\begin{itemize}[leftmargin=1.6em]
\item \textbf{Commutativity}: $E \cup F = F \cup E$, \quad $E \cap F = F \cap E$.
\item \textbf{Associativity}: $E \cup (F \cup G) = (E \cup F) \cup G$, \quad
$E \cap (F \cap G) = (E \cap F) \cap G$.
\item \textbf{Distributivity}: $(E\cup F)\cap G = (E\cap G)\cup(F\cap G)$, \quad
$(E\cap F)\cup G = (E \cup G) \cap (F \cup G)$.
\end{itemize}

\begin{proof}[Proof of one distributive law, as a representative example]
We show $(E \cup F)\cap G = (E \cap G) \cup (F \cap G)$ by mutual inclusion (``element chasing'').

($\subseteq$) Let $\omega \in (E\cup F)\cap G$. Then $\omega \in G$ and ($\omega \in E$ or
$\omega \in F$). If $\omega \in E$, then since also $\omega \in G$, we get $\omega \in E \cap G$,
hence $\omega \in (E\cap G)\cup(F\cap G)$. If instead $\omega \in F$, symmetrically
$\omega \in F \cap G \subset (E\cap G)\cup(F\cap G)$. Either way $\omega$ belongs to the right side.

($\supseteq$) Let $\omega \in (E\cap G)\cup(F\cap G)$. If $\omega \in E \cap G$ then $\omega \in E
\subset E \cup F$ and $\omega \in G$, so $\omega \in (E\cup F)\cap G$. If $\omega \in F \cap G$ the
argument is symmetric. Either way $\omega$ belongs to the left side.

Since each side contains the other, the sets are equal.
\end{proof}

\begin{theorem}[De Morgan's Laws]
For events $E_1,\dots,E_n$ of the same sample space $S$,
\[
\left(\bigcup_{i=1}^n E_i\right)^{\!c} = \bigcap_{i=1}^n E_i^c, \qquad
\left(\bigcap_{i=1}^n E_i\right)^{\!c} = \bigcup_{i=1}^n E_i^c.
\]
\end{theorem}

\begin{proof}
For the first identity: $\omega \in \left(\bigcup_i E_i\right)^c$ $\iff$ $\omega \notin
\bigcup_i E_i$ $\iff$ $\omega \notin E_i$ for every $i$ $\iff$ $\omega \in E_i^c$ for every $i$
$\iff$ $\omega \in \bigcap_i E_i^c$. The chain of iff's establishes set equality. The second
identity follows by applying the first to the events $E_i^c$ (or by an entirely analogous chain of
iff's, replacing every ``for every $i$'' with ``for some $i$'').
\end{proof}

\subsection{The naive definition of probability}
\label{sec:naive-prob}

\begin{definition}[Naive probability]
If a random experiment has a \emph{finite} sample space $S$ in which every sample point is
equally likely, then for $A \subset S$,
\[
\mathbf{P}[A] = \frac{|A|}{|S|} = \frac{\text{number of outcomes favourable to } A}
{\text{total number of outcomes}}.
\]
\end{definition}

This is the definition underlying most of the combinatorial examples of Section~2. It is
intuitive (probability as a \emph{proportion}) but limited: it presupposes finiteness and equal
likelihood, and offers no route to probabilities on infinite or non-uniform sample spaces.

\subsection{The frequentist (relative-frequency) heuristic}

\begin{definition}[Relative-frequency probability]
Suppose a random experiment is repeated $n$ times under (idealized) identical conditions, and let
$n(E)$ denote the number of repetitions in which event $E$ occurs. The relative-frequency notion
of probability sets
\[
\mathbf{P}(E) = \lim_{n \to \infty} \frac{n(E)}{n}.
\]
\end{definition}

This definition, while intuitive, suffers from two logical gaps that the axiomatic theory is
designed to bypass:
\begin{enumerate}
\item There is no a priori guarantee that the sequence $n(E)/n$ converges at all.
\item Even granting convergence, there is no guarantee that the limit is the \emph{same} across
different (infinite) sequences of repetitions of the experiment.
\end{enumerate}

\subsection{The axiomatic (Kolmogorov) definition of probability}

\begin{axiom}[Kolmogorov's axioms]
\label{ax:kolmogorov}
To every event $E$ of the sample space $S$ of a random experiment we associate a real number
$\mathbf{P}(E)$, the \textbf{probability} of $E$, satisfying:
\begin{enumerate}[label=(\arabic*)]
\item $0 \le \mathbf{P}(E) \le 1$ for every event $E \subset S$;
\item $\mathbf{P}(S) = 1$;
\item (Countable additivity) for any sequence $E_1, E_2, \dots$ of pairwise disjoint events (i.e.\
$E_i \cap E_j = \varnothing$ for $i \ne j$),
\[
\mathbf{P}\left(\bigcup_{i=1}^\infty E_i\right) = \sum_{i=1}^\infty \mathbf{P}(E_i).
\]
\end{enumerate}
\end{axiom}

\begin{proposition}
$\mathbf{P}(\varnothing) = 0$.
\end{proposition}
\begin{proof}
Since $S$ and $\varnothing$ are disjoint and $S \cup \varnothing = S$, Axiom~3 (with the sequence
$S, \varnothing, \varnothing, \dots$) gives
\[
\mathbf{P}(S) = \mathbf{P}(S) + \mathbf{P}(\varnothing) + \mathbf{P}(\varnothing) + \cdots.
\]
By Axiom~2, $\mathbf{P}(S)=1$, so $1 = 1 + \sum_{i\ge 1}\mathbf{P}(\varnothing)$. Since each term
$\mathbf{P}(\varnothing) \ge 0$ by Axiom~1, this forces $\mathbf{P}(\varnothing) = 0$.
\end{proof}

\begin{proposition}[Finite additivity]
\label{prop:finite-additivity}
If $E_1, \dots, E_n$ are pairwise disjoint events, then
$\mathbf{P}\left(\bigcup_{i=1}^n E_i\right) = \sum_{i=1}^n \mathbf{P}(E_i)$.
\end{proposition}
\begin{proof}
Extend the finite sequence to an infinite one by setting $E_j := \varnothing$ for all $j > n$. The
extended sequence is still pairwise disjoint ($\varnothing$ is disjoint from every set, including
another copy of $\varnothing$). Axiom~3 gives
\[
\mathbf{P}\left(\bigcup_{i=1}^\infty E_i\right) = \sum_{i=1}^\infty \mathbf{P}(E_i).
\]
But $\bigcup_{i=1}^\infty E_i = \bigcup_{i=1}^n E_i$ (the extra terms are empty), and
$\sum_{i=1}^\infty \mathbf{P}(E_i) = \sum_{i=1}^n \mathbf{P}(E_i) + \sum_{i>n}\mathbf{P}(\varnothing)
= \sum_{i=1}^n \mathbf{P}(E_i)$, using the previous proposition. The claim follows.
\end{proof}

\begin{proposition}[Complement rule]
For any event $E$, $\mathbf{P}(E^c) = 1 - \mathbf{P}(E)$.
\end{proposition}
\begin{proof}
$E$ and $E^c$ are disjoint (by definition of complement) with $E \cup E^c = S$. By finite
additivity and Axiom~2, $1 = \mathbf{P}(S) = \mathbf{P}(E) + \mathbf{P}(E^c)$, whence the claim.
\end{proof}

\begin{proposition}[Monotonicity]
\label{prop:monotone}
If $E \subset F$ (both events of the same sample space), then $\mathbf{P}(E) \le \mathbf{P}(F)$.
\end{proposition}
\begin{proof}
Write $F = E \cup (F\setminus E)$; this is a disjoint union since $E$ and $F \setminus E$ share no
points. By finite additivity, $\mathbf{P}(F) = \mathbf{P}(E) + \mathbf{P}(F\setminus E)$. Since
$\mathbf{P}(F \setminus E) \ge 0$ by Axiom~1, we get $\mathbf{P}(F) \ge \mathbf{P}(E)$.
\end{proof}

\subsection{Inclusion--exclusion}

\begin{theorem}[Inclusion--exclusion, two events]
\label{thm:incl-excl-2}
For any events $E, F$,
\[
\mathbf{P}(E \cup F) = \mathbf{P}(E) + \mathbf{P}(F) - \mathbf{P}(E\cap F).
\]
\end{theorem}
\begin{proof}
The set $E \cup F$ decomposes into the three pairwise disjoint pieces $E\setminus F$,
$F \setminus E$, and $E \cap F$:
\[
E \cup F = (E\setminus F) \cup (F\setminus E) \cup (E\cap F).
\]
Similarly $E = (E\setminus F)\cup(E\cap F)$ (disjoint) and $F=(F\setminus E)\cup(E\cap F)$
(disjoint). Applying finite additivity to each decomposition:
\[
\mathbf{P}(E) = \mathbf{P}(E\setminus F) + \mathbf{P}(E\cap F)
\ \Longrightarrow\ \mathbf{P}(E\setminus F) = \mathbf{P}(E) - \mathbf{P}(E\cap F),
\]
\[
\mathbf{P}(F) = \mathbf{P}(F\setminus E) + \mathbf{P}(E\cap F)
\ \Longrightarrow\ \mathbf{P}(F\setminus E) = \mathbf{P}(F) - \mathbf{P}(E\cap F),
\]
\[
\mathbf{P}(E\cup F) = \mathbf{P}(E\setminus F)+\mathbf{P}(F\setminus E)+\mathbf{P}(E\cap F).
\]
Substituting the first two expressions into the third:
\[
\mathbf{P}(E\cup F) = \big[\mathbf{P}(E)-\mathbf{P}(E\cap F)\big] +
\big[\mathbf{P}(F)-\mathbf{P}(E\cap F)\big] + \mathbf{P}(E\cap F)
= \mathbf{P}(E)+\mathbf{P}(F)-\mathbf{P}(E\cap F). \qedhere
\]
\end{proof}

\begin{theorem}[Inclusion--exclusion, three events]
\[
\mathbf{P}(E\cup F\cup G) = \mathbf{P}(E)+\mathbf{P}(F)+\mathbf{P}(G)
-\mathbf{P}(E\cap F)-\mathbf{P}(E\cap G)-\mathbf{P}(F\cap G)+\mathbf{P}(E\cap F\cap G).
\]
\end{theorem}
\begin{proof}
Apply Theorem~\ref{thm:incl-excl-2} to the two events $E\cup F$ and $G$:
\[
\mathbf{P}(E\cup F\cup G) = \mathbf{P}\big((E\cup F)\cup G\big)
= \mathbf{P}(E\cup F) + \mathbf{P}(G) - \mathbf{P}\big((E\cup F)\cap G\big).
\]
By the distributive law, $(E\cup F)\cap G = (E\cap G)\cup(F\cap G)$, so, applying
Theorem~\ref{thm:incl-excl-2} again to this union:
\[
\mathbf{P}\big((E\cap G)\cup(F\cap G)\big) = \mathbf{P}(E\cap G)+\mathbf{P}(F\cap G)
-\mathbf{P}\big((E\cap G)\cap(F\cap G)\big) = \mathbf{P}(E\cap G)+\mathbf{P}(F\cap G)-\mathbf{P}(E\cap F\cap G),
\]
using $(E\cap G)\cap(F\cap G) = E\cap F\cap G$. Also $\mathbf{P}(E\cup F) = \mathbf{P}(E)+\mathbf{P}(F)-\mathbf{P}(E\cap F)$
by Theorem~\ref{thm:incl-excl-2}. Substituting both expressions above gives the claim.
\end{proof}

\begin{theorem}[General inclusion--exclusion]
\label{thm:incl-excl-n}
For events $E_1,\dots,E_n$ of the same sample space,
\[
\mathbf{P}\left(\bigcup_{i=1}^n E_i\right) = \sum_i \mathbf{P}(E_i) - \sum_{i<j}\mathbf{P}(E_i\cap E_j)
+ \sum_{i<j<k}\mathbf{P}(E_i\cap E_j\cap E_k) - \cdots + (-1)^{n-1}\mathbf{P}\left(\bigcap_{i=1}^n E_i\right).
\]
\end{theorem}
\begin{proof}[Combinatorial proof]
Fix any sample point $\omega \in \bigcup_i E_i$, and suppose $\omega$ belongs to \emph{exactly}
$j \ge 1$ of the $n$ events (relabel so that these are $E_1, \dots, E_j$, for notational
convenience, WLOG). We compute the net number of times $\omega$ is ``counted'' on the right-hand
side.

$\omega$ contributes $1$ to $\mathbf{P}(E_i)$ for each of the $j$ indices $i$ with $\omega \in E_i$
-- i.e.\ $\binom{j}{1}$ contributions at the first-order level. In general, for the $r$-th order
sum $\sum_{i_1<\cdots<i_r}\mathbf{P}(E_{i_1}\cap\cdots\cap E_{i_r})$, $\omega$ contributes to a term
if and only if all $r$ indices lie among the $j$ indices for which $\omega \in E_i$; there are
$\binom{j}{r}$ such terms. Hence the total (signed) contribution of $\omega$ to the right-hand side
is
\[
\sum_{r=1}^{j} (-1)^{r-1}\binom{j}{r}.
\]
We claim this sum equals $1$ for every $j \ge 1$. Indeed, by the binomial theorem
(Theorem~\ref{thm:binomial-thm}) with $x=1,y=-1$:
\[
0 = (1-1)^j = \sum_{r=0}^j \binom{j}{r}(-1)^r = \binom{j}{0} + \sum_{r=1}^j (-1)^r \binom{j}{r}
= 1 - \sum_{r=1}^j (-1)^{r-1}\binom{j}{r},
\]
so $\sum_{r=1}^j (-1)^{r-1}\binom{j}{r} = 1$, as claimed.

Thus every sample point in $\bigcup_i E_i$ contributes net weight exactly $1$ to the right-hand
side, matching its contribution of $1$ to $\mathbf{P}(\bigcup_i E_i)$ (via the naive/relative-
frequency interpretation, or more rigorously, by an entirely analogous inductive argument applying
Theorem~\ref{thm:incl-excl-2} repeatedly, which we omit as it mirrors the $n=3$ case above). A
sample point outside $\bigcup_i E_i$ contributes $0$ to every term on both sides. Hence the two
sides are equal.
\end{proof}

\begin{extra}[Probability as a continuous set function]
This subsection develops material explicitly marked in the source notes as ``further reading, for
the curious'' -- not required for examinations, but presented here with full rigor.

\begin{definition}
A sequence of events $\{E_n\}_{n\in\N}$ is \textbf{increasing} if $E_1 \subset E_2 \subset \cdots$
(i.e.\ $E_i \subset E_{i+1}$ for every $i$), and \textbf{decreasing} if
$E_1 \supset E_2 \supset \cdots$. For an increasing sequence we define
$\displaystyle \lim_{n\to\infty} E_n := \bigcup_{n=1}^\infty E_n$, and for a decreasing sequence,
$\displaystyle \lim_{n\to\infty} E_n := \bigcap_{n=1}^\infty E_n$.
\end{definition}

\begin{theorem}[Continuity of probability]
\label{thm:continuity-prob}
If $\{E_n\}$ is an increasing or decreasing sequence of events, then
\[
\mathbf{P}\left(\lim_{n\to\infty}E_n\right) = \lim_{n\to\infty}\mathbf{P}(E_n).
\]
\end{theorem}

\begin{proof}[Proof (increasing case), via disjointification]
Let $\{E_n\}$ be increasing. Define a new sequence $\{F_n\}$ by
\[
F_1 = E_1, \qquad F_n = E_n \setminus \left(\bigcup_{i=1}^{n-1}E_i\right) = E_n \setminus E_{n-1}
\ \text{for } n \ge 2
\]
(the last equality holds because the sequence is increasing, so $\bigcup_{i=1}^{n-1}E_i=E_{n-1}$).

\emph{Claim 1: the $F_n$ are pairwise disjoint.} Suppose $i<j$. Then $F_j = E_j \setminus
\bigcup_{\ell<j}E_\ell \subset E_j \setminus E_i$ (since $E_i \subset \bigcup_{\ell<j}E_\ell$ for
$i<j$), so $F_j \cap E_i = \varnothing$. Since $F_i \subset E_i$, we get $F_i \cap F_j = \varnothing$.

\emph{Claim 2: $E_n = \bigcup_{i=1}^n F_i$ for every $n$.} We have
$\bigcup_{i=1}^n F_i = F_1 \cup \bigcup_{i=2}^n (E_i \setminus E_{i-1})$. By induction on $n$: true
for $n=1$; assuming $\bigcup_{i=1}^{n-1}F_i = E_{n-1}$, we get
$\bigcup_{i=1}^n F_i = E_{n-1} \cup (E_n \setminus E_{n-1}) = E_{n-1}\cup E_n = E_n$ (since
$E_{n-1}\subset E_n$).

By Claim 2 and finite additivity (Proposition~\ref{prop:finite-additivity}, using Claim 1),
\begin{equation}
\mathbf{P}(E_n) = \mathbf{P}\left(\bigcup_{i=1}^n F_i\right) = \sum_{i=1}^n \mathbf{P}(F_i)
\quad \text{for every } n. \tag{$\ast$}
\end{equation}

\emph{Claim 3: $\bigcup_{n=1}^\infty E_n = \bigcup_{n=1}^\infty F_n$.} If $\omega \in
\bigcup_n E_n$, then $\omega \in E_N$ for some $N$; let $m$ be the \emph{smallest} such index (this
exists by well-ordering of $\N$). Then $\omega \in E_m$ but $\omega \notin E_{m-1}$ (or $m=1$), so
$\omega \in F_m \subset \bigcup_n F_n$. Conversely each $F_n \subset E_n \subset \bigcup_m E_m$, so
$\bigcup_n F_n \subset \bigcup_n E_n$.

By Claim 3, Axiom~3 (countable additivity, valid since the $F_n$ are pairwise disjoint by Claim 1),
and $(\ast)$:
\[
\mathbf{P}\left(\bigcup_{n=1}^\infty E_n\right) = \mathbf{P}\left(\bigcup_{n=1}^\infty F_n\right)
= \sum_{n=1}^\infty \mathbf{P}(F_n) = \lim_{n\to\infty}\sum_{i=1}^n \mathbf{P}(F_i)
= \lim_{n\to\infty}\mathbf{P}(E_n),
\]
which, recalling the definition $\lim_n E_n = \bigcup_n E_n$ for increasing sequences, is exactly
the claim.
\end{proof}

\begin{proof}[Proof (decreasing case), via the increasing case and De Morgan]
If $\{E_n\}$ is decreasing, then $\{E_n^c\}$ is increasing (immediate from $E_n \supset E_{n+1}
\iff E_n^c \subset E_{n+1}^c$). By the increasing case applied to $\{E_n^c\}$:
\[
\lim_{n\to\infty}\mathbf{P}(E_n^c) = \mathbf{P}\left(\bigcup_{n=1}^\infty E_n^c\right)
= \mathbf{P}\left(\Big(\bigcap_{n=1}^\infty E_n\Big)^{\!c}\right)
\]
by De Morgan's law. Using the complement rule on both sides,
\[
\lim_{n\to\infty}\big[1-\mathbf{P}(E_n)\big] = 1 - \mathbf{P}\left(\bigcap_{n=1}^\infty E_n\right)
\ \Longrightarrow\ \lim_{n\to\infty}\mathbf{P}(E_n) = \mathbf{P}\left(\bigcap_{n=1}^\infty E_n\right),
\]
which, since $\lim_n E_n = \bigcap_n E_n$ for decreasing sequences, is the claim.
\end{proof}
\end{extra}

\subsection{Worked examples}

\begin{example}[Statistics course scheduling]
A college has $10$ time slots and offers exactly $3$ statistics courses, each independently and
uniformly randomly assigned one of the $10$ slots (courses may collide). What is the probability
that $2$ or more of the courses share a time slot?

\emph{Solution.} It is easier to compute the complement: ``all three courses get distinct slots.''
The sample space has $10^3 = 1000$ equally likely outcomes (an ordered triple of slot assignments).
The number of outcomes with all-distinct slots is $10\cdot 9\cdot 8 = 720$ (choose a slot for
course 1, then a different slot for course 2, then a different slot for course 3). So
\[
\mathbf{P}(\text{all distinct}) = \frac{720}{1000} = 0.72, \qquad
\mathbf{P}(\text{some collision}) = 1-0.72 = \boxed{0.28}.
\]
\end{example}

\begin{example}[Gimli and Legolas -- linear table]
Twelve people, including Gimli and Legolas, are seated in a row, with every one of the $12!$
arrangements equally likely. Find $\mathbf{P}(\text{Gimli and Legolas are adjacent})$.

\emph{Solution.} Glue Gimli and Legolas into a single ``super-person''; this super-person, together
with the other $10$ people, gives $11$ entities to arrange in $11!$ ways; within the super-person,
Gimli and Legolas can be ordered in $2!$ ways. So the number of favourable arrangements is
$11!\cdot 2!$, and
\[
\mathbf{P}(\text{adjacent}) = \frac{11!\cdot 2!}{12!} = \frac{2}{12} = \boxed{\tfrac{1}{6}}.
\]
\end{example}

\begin{example}[Gimli and Legolas -- circular table]
The same twelve people are now seated around a circular table with unlabeled chairs, but
left/right position relative to others still matters (rotations of a seating are considered the
\emph{same} arrangement; reflections are considered \emph{different}). Find
$\mathbf{P}(\text{Gimli and Legolas adjacent})$.

\emph{Solution.} The number of distinct circular arrangements of $n$ distinct people (identifying
rotations) is $(n-1)!$: fix one person's seat as a reference (this uses up the rotational
symmetry), and arrange the remaining $n-1$ people in the remaining seats in $(n-1)!$ ways. So the
total sample space has $11! $ equally likely arrangements.

For the favourable count, glue Gimli and Legolas into a super-person; now $11$ entities are
arranged around the circle, giving $(11-1)! = 10!$ circular arrangements, times $2!$ internal
orders of the super-person. So
\[
\mathbf{P}(\text{adjacent}) = \frac{10!\cdot 2!}{11!} = \frac{2}{11}.
\]
\emph{(Note this probability is slightly higher than in the linear case, $2/11 > 1/6$, since on a
circle every seat has two neighbours, including the two ``end'' seats of the linear arrangement,
which only have one neighbour each.)}
\end{example}

\begin{example}[The matching / derangement problem, and the coat-check problem]
\label{ex:derangement}
A friend has $n$ cards numbered $1,\dots,n$, placed face-down on a table in an order unknown to
you. You guess the number on each card, in order, from left to right. You \emph{win} if at least
one guess matches. Equivalently (the ``coat check'' formulation): $n$ party-goers each receive a
uniformly random permutation of $n$ coats; what is the probability that \emph{no one} receives
their own coat back (a \textbf{derangement})?

\emph{Solution.} Let $A_i$ be the event that the $i$-th person receives their own coat (equivalently,
guess $i$ is correct). We want $\mathbf{P}\!\left[\bigcap_{i=1}^n A_i^c\right] = 1 - \mathbf{P}\!\left[\bigcup_{i=1}^n A_i\right]$.

For any distinct indices $i_1,\dots,i_k \in \{1,\dots,n\}$, conditioning on $k$ specific people
getting their own coats leaves the remaining $n-k$ coats distributed uniformly at random among the
remaining $n-k$ people, so
\[
\mathbf{P}[A_{i_1}\cap \cdots \cap A_{i_k}] = \frac{(n-k)!}{n!}.
\]
By Theorem~\ref{thm:incl-excl-n}, noting there are $\dbinom{n}{k}$ terms of order $k$
(each with the same value $(n-k)!/n!$ by symmetry):
\[
\mathbf{P}\!\left[\bigcup_{i=1}^n A_i\right] = \sum_{k=1}^n (-1)^{k-1}\binom{n}{k}\frac{(n-k)!}{n!}.
\]
Hence
\[
\mathbf{P}[\text{no matches}] = 1 - \sum_{k=1}^n (-1)^{k-1}\binom{n}{k}\frac{(n-k)!}{n!}
= \sum_{k=0}^n (-1)^k \binom{n}{k}\frac{(n-k)!}{n!} = \sum_{k=0}^n \frac{(-1)^k}{k!},
\]
using $\binom{n}{k}\frac{(n-k)!}{n!} = \frac{1}{k!}$. So
\[
\boxed{\mathbf{P}[\text{no matches}] = \sum_{k=0}^n \frac{(-1)^k}{k!}}
\qquad\text{and}\qquad
\mathbf{P}[\text{you win}] = 1 - \sum_{k=0}^n \frac{(-1)^k}{k!}.
\]
Since $\sum_{k=0}^\infty (-1)^k/k! = e^{-1}$, for large $n$,
$\mathbf{P}[\text{no matches}] \approx e^{-1} \approx 0.3679$, so
$\mathbf{P}[\text{you win}] \to 1-e^{-1} \approx 0.6321$, remarkably independent of $n$ (for large
$n$).
\end{example}

\begin{example}[Bayes/partition problem: the missing plane]
A plane is missing and, a priori, equally likely to have gone down in any of $3$ regions
($\mathbf{P}(R_i)=1/3$ for $i=1,2,3$). Let $1-\beta_i$ denote the probability that a search of
region $i$ finds the plane, \emph{given} the plane is actually in region $i$. Given that a search
of region $1$ was unsuccessful, find $\mathbf{P}(R_i \mid F_1^c)$ for each $i$, where $F_1^c$
denotes ``the search of region 1 failed.''

\emph{Solution.} If the plane is not in region $1$, a search of region $1$ is certain to fail:
$\mathbf{P}(F_1^c \mid R_2) = \mathbf{P}(F_1^c\mid R_3) = 1$. If the plane is in region 1,
$\mathbf{P}(F_1^c\mid R_1) = \beta_1$ by definition. By the law of total probability (established
formally in Section~4),
\[
\mathbf{P}(F_1^c) = \beta_1\cdot\tfrac13 + 1\cdot \tfrac13 + 1\cdot\tfrac13 = \frac{\beta_1+2}{3}.
\]
By Bayes' formula,
\[
\mathbf{P}(R_1\mid F_1^c) = \frac{\beta_1 \cdot \tfrac13}{(\beta_1+2)/3} = \frac{\beta_1}{\beta_1+2},
\qquad
\mathbf{P}(R_2\mid F_1^c) = \mathbf{P}(R_3\mid F_1^c) = \frac{1\cdot \tfrac13}{(\beta_1+2)/3} =
\frac{1}{\beta_1+2}.
\]
As a sanity check, these three probabilities sum to $\dfrac{\beta_1+1+1}{\beta_1+2}=1$. Note that
$\mathbf{P}(R_1\mid F_1^c) < \mathbf{P}(R_1) = 1/3$ (an unsuccessful search of region 1 makes it
\emph{less} likely the plane is there, and correspondingly more likely it is in region 2 or 3) --
unless $\beta_1=1$ (search is useless), in which case the failed search carries no information and
$\mathbf{P}(R_i\mid F_1^c) = 1/3$ for all $i$.
\end{example}

\begin{example}[Bayes' formula: a diagnostic test]
A disease test is $95\%$ sensitive ($\mathbf{P}[+\mid \text{disease}]=0.95$) and has a $1\%$
false-positive rate ($\mathbf{P}[+\mid \text{no disease}] = 0.01$). The disease prevalence is
$0.5\%$. Given a positive test result, find the probability the person actually has the disease.

\emph{Solution.} By the law of total probability,
\[
\mathbf{P}[+] = \mathbf{P}[+\mid D]\mathbf{P}[D] + \mathbf{P}[+\mid D^c]\mathbf{P}[D^c]
= (0.95)(0.005) + (0.01)(0.995) = 0.00475 + 0.00995 = 0.0147.
\]
By Bayes' formula,
\[
\mathbf{P}[D\mid +] = \frac{\mathbf{P}[+\mid D]\mathbf{P}[D]}{\mathbf{P}[+]}
= \frac{0.00475}{0.0147} \approx \boxed{0.3231}.
\]
Despite the test's apparent accuracy, fewer than a third of positive results are true positives --
a consequence of the low base rate (prevalence). This is the classic caution against confusing
$\mathbf{P}[+\mid D]$ with $\mathbf{P}[D\mid +]$.
\end{example}

\newpage
% =========================================================================
\section{Conditional Probability and Independence}
% =========================================================================

\subsection{Definition and basic examples}

\begin{definition}
Given events $E,F$ with $\mathbf{P}[F]>0$, the \textbf{conditional probability} of $E$ given $F$
is
\[
\mathbf{P}[E\mid F] := \frac{\mathbf{P}[E\cap F]}{\mathbf{P}[F]}.
\]
\end{definition}

\begin{intuition}
Knowing that $F$ has occurred restricts attention to the ``reduced sample space'' $F$: outcomes
outside $F$ are now known to be impossible. Within this reduced space, the (renormalized)
likelihood of $E$ occurring is the fraction of $F$ that also lies in $E$, i.e.\
$|E\cap F|/|F|$ in the naive/finite-uniform case -- exactly $\mathbf{P}[E\cap F]/\mathbf{P}[F]$.
\end{intuition}

\begin{proposition}
For fixed $F$ with $\mathbf{P}[F]>0$, the set function $f(A) := \mathbf{P}[A\mid F]$, defined for
events $A \subset S$, is itself a legitimate probability function, i.e.\ satisfies
Axiom~\ref{ax:kolmogorov}.
\end{proposition}
\begin{proof}
(1) $f(A) = \mathbf{P}[A\cap F]/\mathbf{P}[F] \ge 0$ since numerator and denominator are
non-negative, and $f(A) \le 1$ since $\mathbf{P}[A\cap F]\le \mathbf{P}[F]$ by monotonicity
(Proposition~\ref{prop:monotone}, as $A\cap F \subset F$).
(2) $f(S) = \mathbf{P}[S\cap F]/\mathbf{P}[F] = \mathbf{P}[F]/\mathbf{P}[F]=1$.
(3) If $A_1,A_2,\dots$ are pairwise disjoint, so are $A_1\cap F, A_2\cap F,\dots$, and
$\left(\bigcup_i A_i\right)\cap F = \bigcup_i (A_i \cap F)$ by distributivity, so
\[
f\left(\bigcup_i A_i\right) = \frac{\mathbf{P}\left[\bigcup_i(A_i\cap F)\right]}{\mathbf{P}[F]}
= \frac{\sum_i \mathbf{P}[A_i\cap F]}{\mathbf{P}[F]} = \sum_i f(A_i). \qedhere
\]
\end{proof}

\subsection{The multiplication rule}

\begin{theorem}[Multiplication rule]
For events $E_1,\dots,E_n$ with $\mathbf{P}[E_1\cap\cdots\cap E_{n-1}]>0$,
\[
\mathbf{P}\!\left[\bigcap_{i=1}^n E_i\right] = \mathbf{P}[E_1]\,\mathbf{P}[E_2\mid E_1]\,
\mathbf{P}[E_3\mid E_1\cap E_2]\cdots \mathbf{P}\!\left[E_n \,\middle|\, \bigcap_{i=1}^{n-1}E_i\right].
\]
\end{theorem}
\begin{proof}
By definition, $\mathbf{P}[E_2\mid E_1]\mathbf{P}[E_1] = \mathbf{P}[E_1\cap E_2]$. Inductively,
$\mathbf{P}\!\left[E_k \mid \bigcap_{i<k}E_i\right]\cdot \mathbf{P}\!\left[\bigcap_{i<k}E_i\right]
= \mathbf{P}\!\left[\bigcap_{i\le k}E_i\right]$ directly from the definition of conditional
probability (with $E=E_k$, $F=\bigcap_{i<k}E_i$). Chaining these equalities from $k=n$ down to
$k=2$ telescopes to the stated product.
\end{proof}

\subsection{The law of total probability and Bayes' formula}

\begin{theorem}[Law of total probability -- two-event form]
For events $E,F$ with $0<\mathbf{P}[F]<1$,
\[
\mathbf{P}[E] = \mathbf{P}[E\mid F]\,\mathbf{P}[F] + \mathbf{P}[E\mid F^c]\,\mathbf{P}[F^c].
\]
\end{theorem}
\begin{proof}
$E = (E\cap F)\cup(E\cap F^c)$ is a disjoint union, so by finite additivity,
$\mathbf{P}[E]=\mathbf{P}[E\cap F]+\mathbf{P}[E\cap F^c]
= \mathbf{P}[E\mid F]\mathbf{P}[F] + \mathbf{P}[E\mid F^c]\mathbf{P}[F^c]$, using the definition of
conditional probability twice (in the form $\mathbf{P}[E\cap F]=\mathbf{P}[E\mid F]\mathbf{P}[F]$).
\end{proof}

\begin{theorem}[Law of total probability -- partition form]
Let $E_1,\dots,E_n$ be a \textbf{partition} of the sample space $\Omega$, i.e.\
$\bigcup_i E_i = \Omega$ and $E_i \cap E_j = \varnothing$ for $i \ne j$, each with
$\mathbf{P}[E_i]>0$. Then for any event $F$,
\[
\mathbf{P}[F] = \sum_{i=1}^n \mathbf{P}[F\mid E_i]\,\mathbf{P}[E_i].
\]
\end{theorem}
\begin{proof}
$F = F \cap \Omega = F \cap \bigcup_i E_i = \bigcup_i (F\cap E_i)$ (distributivity), a disjoint
union since the $F\cap E_i$ inherit disjointness from the $E_i$. By finite additivity,
$\mathbf{P}[F] = \sum_i \mathbf{P}[F\cap E_i] = \sum_i \mathbf{P}[F\mid E_i]\mathbf{P}[E_i]$.
\end{proof}

\begin{theorem}[Bayes' formula, general form]
\label{thm:bayes}
With $E_1,\dots,E_n$ a partition of $\Omega$ as above, and $F$ an event with $\mathbf{P}[F]>0$,
\[
\mathbf{P}[E_i \mid F] = \frac{\mathbf{P}[F\mid E_i]\,\mathbf{P}[E_i]}
{\sum_{j=1}^n \mathbf{P}[F\mid E_j]\,\mathbf{P}[E_j]}.
\]
\end{theorem}
\begin{proof}
By definition, $\mathbf{P}[E_i\mid F] = \mathbf{P}[F\cap E_i]/\mathbf{P}[F] =
\mathbf{P}[F\mid E_i]\mathbf{P}[E_i]/\mathbf{P}[F]$; now substitute the law of total probability
for the denominator $\mathbf{P}[F]$.
\end{proof}

\subsection{Worked examples}

\begin{example}[Bridge hands]
A standard $52$-card deck is dealt equally to North, South, East, West (13 cards each, uniformly
at random). Given that North and South together hold $8$ spades, find $\mathbf{P}(\text{East holds
exactly 3 spades})$.

\emph{Solution.} There are $13$ spades total; if North+South (26 cards) hold $8$ of them, the
remaining $26$ cards (held jointly by East and West) contain the remaining $13-8=5$ spades. Given
this, the $26$ non-(North/South) cards are dealt uniformly at random into East's $13$ and West's
$13$; this is exactly a hypergeometric sampling situation: choosing $13$ of $26$ cards (5 spades,
21 non-spades) for East. Hence
\[
\mathbf{P}(\text{East has exactly 3 spades}) = \frac{\dbinom{5}{3}\dbinom{21}{10}}{\dbinom{26}{13}}
= \frac{10\cdot 352716}{10400600} \approx \boxed{0.3391}.
\]
\end{example}

\begin{example}[Aces in four bridge hands]
The $52$ cards are dealt into $4$ labeled piles of $13$ each. Find $\mathbf{P}(\text{each pile
contains exactly one ace})$.

\emph{Solution 1 (multiplication rule).} Deal the four aces one at a time (conceptually), tracking
which pile each lands in, along with the other $48$ cards filling in the remaining $12$ slots per
pile. The first ace can go anywhere. Given the first ace's pile, the second ace lands in a
\emph{different} pile with probability $39/51$ (there are $51$ remaining card-slots overall, of
which $12$ remain in the first ace's pile and $39$ lie in the other three piles). Given the first
two aces are in different piles, the third ace avoids both with probability $26/50$ (of the
remaining $50$ slots, $26$ lie in the two untouched piles). The fourth ace must then land in the
one remaining (ace-free) pile, with probability $13/49$. By the multiplication rule,
\[
\mathbf{P} = \frac{39}{51}\cdot\frac{26}{50}\cdot\frac{13}{49}
= \frac{39\cdot26\cdot13}{51\cdot50\cdot49} = \frac{13182}{124950} \approx \boxed{0.1055}.
\]

\emph{Solution 2 (direct counting).} The total number of ways to deal $52$ cards into $4$ labeled
piles of $13$ is the multinomial coefficient $\dbinom{52}{13,13,13,13} = \dfrac{52!}{(13!)^4}$.
Favourable deals: first assign the $4$ distinct aces to the $4$ piles, one each, in $4!$ ways; then
distribute the remaining $48$ cards into the remaining $12$ slots of each pile, in
$\dbinom{48}{12,12,12,12}=\dfrac{48!}{(12!)^4}$ ways. So
\[
\mathbf{P} = \frac{4!\cdot 48!/(12!)^4}{52!/(13!)^4}
= 4!\left(\frac{13!}{12!}\right)^{\!4}\frac{48!}{52!}
= \frac{24\cdot 13^4}{52\cdot51\cdot50\cdot49} = \frac{685464}{6497400}\approx 0.1055,
\]
agreeing exactly with Solution 1.
\end{example}

\begin{example}[The urn model without replacement, and exchangeability]
\label{ex:urn}
An urn contains $b$ black and $w$ white balls; balls are drawn sequentially, without replacement.

\textbf{(a)} $\mathbf{P}(\text{1st black, 2nd white})$ vs.\ $\mathbf{P}(\text{1st white, 2nd
black})$:
\[
\mathbf{P}(B_1 W_2) = \frac{b}{b+w}\cdot\frac{w}{b+w-1}, \qquad
\mathbf{P}(W_1 B_2) = \frac{w}{b+w}\cdot\frac{b}{b+w-1}.
\]
These are \emph{equal} -- both equal $\dfrac{bw}{(b+w)(b+w-1)}$ -- even though the events describe
opposite orders.

\textbf{(b) General exchangeability.} For any $k \le b+w$ and any specific sequence of colours
$C_1,\dots,C_k$ (each black or white) with exactly $\ell$ black and $k-\ell$ white,
\[
\mathbf{P}[\text{draws } 1,\dots,k \text{ have colours } C_1,\dots,C_k]
= \frac{b(b-1)\cdots(b-\ell+1)\;w(w-1)\cdots(w-k+\ell+1)}{(b+w)(b+w-1)\cdots(b+w-k+1)},
\]
\emph{regardless of the particular order in which the $\ell$ blacks and $k-\ell$ whites appear.}

\emph{Proof.} Every time a ball is drawn, its (conditional) probability of being black equals
(current number of black balls)/(current total number of balls). The denominator, at the $i$-th
draw, is $b+w-i+1$ \emph{regardless of the colours drawn previously}; multiplying over
$i=1,\dots,k$ gives the denominator $(b+w)(b+w-1)\cdots(b+w-k+1)$. For the numerator: each time a
black ball is drawn, the number of black balls available decreases by one from its current value,
regardless of when the black draws occur relative to white draws; so if black balls are drawn at
$\ell$ of the $k$ draw-times (in any pattern), the successive black-counts realized are
$b, b-1,\dots,b-\ell+1$ (in that order, since each black draw reduces the black-count by exactly
one from wherever it currently stands) -- and likewise the white contributions are
$w,w-1,\dots,w-k+\ell+1$. The product of all these numerators, by the multiplication rule, is as
stated, independent of the interleaving pattern. This property is called \textbf{exchangeability}.
$\blacksquare$

\textbf{(c) Given $k$ black balls drawn among the first $n$: probability the first ball is
black.} By part (b), \emph{every} sequence of $n$ colours containing exactly $k$ blacks has the
\emph{same} probability; there are $\dbinom{n}{k}$ such sequences (choose which of the $n$ draw-
positions are black), so, conditionally on ``exactly $k$ blacks among the first $n$ draws,'' each
of these $\dbinom{n}{k}$ sequences is equally likely. The number of these sequences with the first
draw black is $\dbinom{n-1}{k-1}$ (the remaining $k-1$ blacks are placed among the other $n-1$
positions). Hence
\[
\mathbf{P}(\text{1st draw black} \mid k \text{ blacks among first } n)
= \frac{\dbinom{n-1}{k-1}}{\dbinom{n}{k}} = \frac{k}{n}.
\]
\emph{(This last simplification: $\dbinom{n-1}{k-1}/\dbinom{n}{k} =
\dfrac{(n-1)!/[(k-1)!(n-k)!]}{n!/[k!(n-k)!]} = \dfrac{(n-1)!\,k!}{(k-1)!\,n!}=\dfrac{k}{n}$.)}
\end{example}

\subsection{Independence}

\begin{definition}
Events $E,F$ (with $\mathbf{P}[F]>0$) are \textbf{independent} if $\mathbf{P}[E\mid F]=\mathbf{P}[E]$,
equivalently (multiplying both sides by $\mathbf{P}[F]$),
\[
\mathbf{P}[E\cap F] = \mathbf{P}[E]\,\mathbf{P}[F].
\]
\end{definition}

\begin{proposition}
If $E$ and $F$ are independent, so are $E$ and $F^c$.
\end{proposition}
\begin{proof}
$\mathbf{P}[E\cap F^c] = \mathbf{P}[E] - \mathbf{P}[E\cap F]$ (from $E = (E\cap F)\cup(E\cap F^c)$,
disjoint union, and finite additivity) $= \mathbf{P}[E] - \mathbf{P}[E]\mathbf{P}[F]$ (using
independence of $E,F$) $= \mathbf{P}[E](1-\mathbf{P}[F]) = \mathbf{P}[E]\mathbf{P}[F^c]$.
\end{proof}

\begin{definition}
A finite collection of events $E_1,\dots,E_r$ is \textbf{(mutually) independent} if, for
\emph{every} nonempty subset of indices $\{i_1,\dots,i_s\}\subset\{1,\dots,r\}$,
\[
\mathbf{P}\!\left[\bigcap_{j=1}^s E_{i_j}\right] = \prod_{j=1}^s \mathbf{P}[E_{i_j}].
\]
An infinite collection is independent if every finite sub-collection is independent.
\end{definition}

\begin{example}[Pairwise independence does not imply mutual independence]
\label{ex:pairwise-not-mutual}
Roll two fair, independent dice. Let
\[
A = \{\text{die 1 is even}\}, \quad B=\{\text{die 2 is even}\}, \quad
C=\{\text{both dice show the same parity}\}.
\]
Each of $A,B,C$ has probability $1/2$ (by symmetry: $C$ occurs for $18$ of the $36$ equally likely
outcomes -- $9$ with both even, $9$ with both odd).

\emph{Pairwise independence.} $\mathbf{P}[A\cap B]$ = both even = $\tfrac12\cdot\tfrac12=\tfrac14 =
\mathbf{P}[A]\mathbf{P}[B]$: independent. $\mathbf{P}[A\cap C]$ = die 1 even \emph{and} same parity
= die 1 and die 2 both even = $\tfrac14 = \mathbf{P}[A]\mathbf{P}[C]$: independent. By symmetry
$B,C$ are independent too.

\emph{Failure of mutual independence.} $\mathbf{P}[A\cap B\cap C]$ = both dice even (which
automatically satisfies $C$) $=\tfrac14$. But $\mathbf{P}[A]\mathbf{P}[B]\mathbf{P}[C] =
\tfrac12\cdot\tfrac12\cdot\tfrac12=\tfrac18 \ne \tfrac14$.

So $A,B,C$ are pairwise independent but \emph{not} mutually independent: knowing any two of $A,B,C$
in fact fully determines the third (e.g.\ $A$ and $C$ together force $B$), so the triple carries
strictly less ``independent information'' than three genuinely unrelated coin flips would.
\end{example}

\subsection{Comparison table: conditional probability concepts}

\begin{center}
\renewcommand{\arraystretch}{1.5}
\begin{longtable}{p{3.3cm} p{5.7cm} p{5.7cm}}
\toprule
\textbf{Concept} & \textbf{Definition} & \textbf{Key identity} \\
\midrule
Conditional probability & Belief in $E$ updated by knowing $F$ occurred & $\mathbf{P}[E\mid F]=\mathbf{P}[E\cap F]/\mathbf{P}[F]$ \\
Independence (2 events) & $F$ carries no information about $E$ & $\mathbf{P}[E\cap F]=\mathbf{P}[E]\mathbf{P}[F]$ \\
Pairwise independence ($n$ events) & every \emph{pair} multiplies & $\mathbf{P}[E_i\cap E_j]=\mathbf{P}[E_i]\mathbf{P}[E_j]$, all $i\ne j$ \\
Mutual independence ($n$ events) & every \emph{subcollection} multiplies & $\mathbf{P}\big[\bigcap_{j\in S}E_j\big]=\prod_{j\in S}\mathbf{P}[E_j]$, all $S$ \\
Law of total probability & Decompose $F$ over a partition & $\mathbf{P}[F]=\sum_i \mathbf{P}[F\mid E_i]\mathbf{P}[E_i]$ \\
Bayes' formula & Invert a conditional probability & $\mathbf{P}[E_i\mid F] = \mathbf{P}[F\mid E_i]\mathbf{P}[E_i]/\mathbf{P}[F]$ \\
\bottomrule
\end{longtable}
\end{center}

\newpage
% =========================================================================
\section{Discrete Random Variables}
% =========================================================================

\subsection{Definition, pmf, and cdf}

\begin{definition}
A \textbf{random variable} is a (real-valued, or more generally $\R^n$-valued) function
$X : S \to \R$ defined on the sample space of a random experiment.
\end{definition}

\begin{definition}
A random variable $X$ is \textbf{discrete} if it takes at most countably many values.
The \textbf{probability mass function (pmf)} of a discrete $X$ is
$p(x) := \mathbf{P}[X=x]$ for $x \in \R$; necessarily $p(x)>0$ for at most countably many $x$, and
$\sum_{i} p(x_i) = 1$ where $x_1,x_2,\dots$ enumerate the values with $p(x_i)>0$.
\end{definition}

\begin{definition}
The \textbf{cumulative distribution function (cdf)} of any random variable $X$ is
$F(a) := \mathbf{P}[X\le a]$ for $a \in \R$. For discrete $X$, $F(a) = \sum_{x\le a}p(x)$.
\end{definition}

\begin{proposition}[Properties of the cdf]
\label{prop:cdf-properties}
For any random variable $X$ with cdf $F$:
\begin{enumerate}[label=(\arabic*)]
\item (Monotonicity) $x_1 \le x_2 \implies F(x_1)\le F(x_2)$.
\item (Right-continuity) $\lim_{y\to x^+}F(y) = F(x)$ for every $x$.
\item (Limits at infinity) $\lim_{x\to\infty}F(x)=1$, $\lim_{x\to-\infty}F(x)=0$.
\end{enumerate}
\end{proposition}
\begin{proof}
\textbf{(1)} $\{X\le x_1\}\subset \{X\le x_2\}$ as sets (any $\omega$ with $X(\omega)\le x_1$ also
satisfies $X(\omega)\le x_2$), so by monotonicity of probability (Proposition~\ref{prop:monotone}),
$F(x_1)=\mathbf{P}[X\le x_1]\le \mathbf{P}[X\le x_2]=F(x_2)$.

\textbf{(2)} Take any sequence $x_n \downarrow x$ (decreasing to $x$), and set $A_n := \{X\le x_n\}$.
Since $x_n$ is decreasing, $\{A_n\}$ is a decreasing sequence of events, with
$\bigcap_n A_n = \{X\le x\}$ (an outcome satisfies $X(\omega)\le x_n$ for \emph{every} $n$ if and
only if $X(\omega)\le x$, since $x_n \downarrow x$). By continuity of probability
(Theorem~\ref{thm:continuity-prob}), $F(x_n)=\mathbf{P}[A_n]\to \mathbf{P}[\bigcap_n A_n]=F(x)$.
Since this holds for \emph{every} sequence $x_n\downarrow x$, we conclude $\lim_{y\to x^+}F(y)=F(x)$.

\textbf{(3)} Apply Theorem~\ref{thm:continuity-prob} to the increasing sequence
$\{X\le n\}_{n\in\N}$, whose union is all of $S$ (every real-valued $X(\omega)$ is $\le$ some
integer), giving $F(n)\to \mathbf{P}(S)=1$; and to the decreasing sequence $\{X\le -n\}_{n\in\N}$,
whose intersection is $\varnothing$, giving $F(-n)\to \mathbf{P}(\varnothing)=0$.
\end{proof}

\begin{remark}
For a discrete random variable, $F$ is a right-continuous \textbf{step function}, with jumps of
size $p(x)$ exactly at the points $x$ where $p(x)>0$. For a continuous random variable (Section~8),
by contrast, $F$ is continuous everywhere and $\mathbf{P}[X=x]=0$ for every single point $x$.
\end{remark}

\subsection{Expectation and the law of the unconscious statistician (LOTUS)}

\begin{definition}
The \textbf{expectation} of a discrete random variable $X$ with pmf $p$ is
\[
\E{X} = \sum_{x\,:\,p(x)>0} x\,p(x).
\]
\end{definition}

\begin{theorem}[LOTUS]
\label{thm:lotus}
Let $X$ be discrete, taking values $x_1,x_2,\dots$, with pmf $p$, and let $g:\R\to\R$. Then
$Y := g(X)$ is also a discrete random variable, and
\[
\E{g(X)} = \sum_{i} g(x_i)\,p(x_i).
\]
\end{theorem}
\begin{proof}
Let $Y$ take values in $\{y_1,y_2,\dots\}$ (a countable set, since $Y$ is a function of the
countably-valued $X$). For each $j$, define $S_j := \{x_i : g(x_i)=y_j\}$. The sets $\{S_j\}$ form
a \textbf{partition} of $\{x_1,x_2,\dots\}$: they are disjoint (an $x_i$ maps to a unique $g(x_i)$,
hence lies in exactly one $S_j$), and their union is everything (every $x_i$ lies in $S_j$ for
$j$ such that $y_j = g(x_i)$).

Since $Y=y_j$ if and only if $X=x_i$ for some $x_i \in S_j$, the events $\{X=x_i\}$, $x_i \in S_j$,
are disjoint with union $\{Y=y_j\}$, so by finite (or countable) additivity,
\[
\mathbf{P}[Y=y_j] = \sum_{x_i \in S_j} \mathbf{P}[X=x_i] = \sum_{x_i\in S_j}p(x_i).
\]
Now compute, using the partition $\{S_j\}$ to reorganize the sum (valid since all terms are
non-negative, or more generally the series converges absolutely -- an implicit standing
assumption whenever we write $\E{X}$):
\[
\sum_i g(x_i)p(x_i) = \sum_j \sum_{x_i \in S_j} g(x_i)p(x_i)
= \sum_j \sum_{x_i\in S_j} y_j\, p(x_i)
= \sum_j y_j \sum_{x_i \in S_j}p(x_i) = \sum_j y_j\, \mathbf{P}[Y=y_j] = \E{Y}. \qedhere
\]
\end{proof}

\begin{remark}
The value of LOTUS is practical: to compute $\E{g(X)}$ one need not first derive the pmf of the
transformed variable $Y=g(X)$; the pmf of $X$ together with $g$ suffices.
\end{remark}

\begin{corollary}[Linearity of expectation, discrete scalar case]
For constants $a,b\in\R$, $\E{aX+b} = a\E{X}+b$.
\end{corollary}
\begin{proof}
Apply LOTUS with $g(x)=ax+b$: $\E{aX+b} = \sum_i (ax_i+b)p(x_i) = a\sum_i x_i p(x_i) + b\sum_i p(x_i)
= a\E{X}+b\cdot 1$.
\end{proof}

\subsection{The discrete distributions catalogue}

We now derive, from first principles, the pmf, mean, and variance of each of the standard discrete
distributions appearing in the notes.

\subsubsection{Bernoulli$(p)$}

A single trial with $\mathbf{P}[X=1]=p$ (``success''), $\mathbf{P}[X=0]=1-p$ (``failure'').
\[
\E{X} = 1\cdot p + 0\cdot(1-p) = p, \qquad
\E{X^2} = 1^2\cdot p+0^2\cdot(1-p) = p,\qquad
\Var{X} = p-p^2 = p(1-p).
\]

\subsubsection{Binomial$(n,p)$}

$n$ independent Bernoulli$(p)$ trials; $X=$ total number of successes.

\begin{theorem}
$\displaystyle \mathbf{P}[X=k] = \binom{n}{k}p^k(1-p)^{n-k}$, $k=0,1,\dots,n$.
\end{theorem}
\begin{proof}
Fix $k$. Any \emph{specific} sequence of $n$ trial-outcomes with exactly $k$ successes and $n-k$
failures has probability $p^k(1-p)^{n-k}$ by independence (multiplication rule for independent
events). The number of such sequences is the number of ways to choose \emph{which} $k$ of the $n$
trials are the successes, i.e.\ $\dbinom{n}{k}$. Since these $\dbinom{n}{k}$ sequences are disjoint
events whose union is $\{X=k\}$, finite additivity gives the stated pmf.
\end{proof}

\begin{theorem}
$\E{X}=np$ and $\Var{X}=np(1-p)$.
\end{theorem}
\begin{proof}
\textbf{Mean.} Write $X = \sum_{i=1}^n X_i$ where $X_i = \mathbf{1}_{\{\text{trial } i \text{ is a
success}\}}$ is Bernoulli$(p)$. By linearity of expectation (extended by induction to $n$
summands, from the two-summand case proved in Section~6),
$\E{X} = \sum_i \E{X_i} = \sum_i p = np$.

\textbf{Variance,} via the factorial moment $\E{X(X-1)}$ (this route is self-contained and avoids
needing to prove that variances of independent variables add):
\[
\E{X(X-1)} = \sum_{k=0}^n k(k-1)\binom{n}{k}p^k(1-p)^{n-k}.
\]
The $k=0,1$ terms vanish. For $k\ge 2$, use
$k(k-1)\dbinom{n}{k} = k(k-1)\dfrac{n!}{k!(n-k)!} = n(n-1)\dfrac{(n-2)!}{(k-2)!(n-k)!}
= n(n-1)\dbinom{n-2}{k-2}$. So
\[
\E{X(X-1)} = n(n-1)p^2\sum_{k=2}^n \binom{n-2}{k-2}p^{k-2}(1-p)^{n-k}
= n(n-1)p^2 \sum_{j=0}^{n-2}\binom{n-2}{j}p^j(1-p)^{(n-2)-j} = n(n-1)p^2,
\]
where the last sum is $1$ by the binomial theorem (it is the total probability of a Binomial$(n-2,p)$
pmf). Hence
\[
\E{X^2} = \E{X(X-1)}+\E{X} = n(n-1)p^2+np,
\]
and
\[
\Var{X} = \E{X^2}-(\E{X})^2 = n(n-1)p^2+np - n^2p^2 = np - np^2 = np(1-p). \qedhere
\]
\end{proof}

\subsubsection{Geometric$(p)$}

$X=$ number of independent Bernoulli$(p)$ trials up to and including the first success.

\begin{theorem}
$\mathbf{P}[X=k]=(1-p)^{k-1}p$, $k=1,2,\dots$; $\E{X}=1/p$; $\Var{X}=(1-p)/p^2$.
\end{theorem}
\begin{proof}
\textbf{PMF.} $\{X=k\}$ = ``first $k-1$ trials fail, $k$-th succeeds''; by independence, this has
probability $(1-p)^{k-1}p$.

\textbf{Mean.} Write $q=1-p$. Differentiate the geometric series $\sum_{k=0}^\infty q^k = 1/(1-q)$
term-by-term with respect to $q$ (valid for $|q|<1$, a standard fact about power series within
their radius of convergence):
\[
\sum_{k=1}^\infty k q^{k-1} = \frac{1}{(1-q)^2}.
\]
Then $\E{X} = \sum_{k=1}^\infty k q^{k-1}p = p\cdot\dfrac{1}{(1-q)^2} = \dfrac{p}{p^2}=\dfrac1p$.

\textbf{Variance.} Differentiate once more: $\sum_{k=2}^\infty k(k-1)q^{k-2} = \dfrac{2}{(1-q)^3}$,
so $\sum_{k=1}^\infty k(k-1)q^{k-1} = q\cdot\dfrac{2}{(1-q)^3} = \dfrac{2q}{p^3}$. Then
\[
\E{X(X-1)} = p\sum_{k=1}^\infty k(k-1)q^{k-1} = p\cdot\frac{2q}{p^3} = \frac{2q}{p^2},
\]
\[
\E{X^2} = \E{X(X-1)}+\E{X} = \frac{2q}{p^2}+\frac1p = \frac{2q+p}{p^2} = \frac{2-p}{p^2}
\quad(\text{using } q=1-p),
\]
\[
\Var{X} = \E{X^2}-(\E{X})^2 = \frac{2-p}{p^2}-\frac{1}{p^2} = \frac{1-p}{p^2}. \qedhere
\]
\end{proof}

\subsubsection{Negative Binomial$(k,p)$}

$Y=$ number of independent Bernoulli$(p)$ trials up to and including the $k$-th success.

\begin{theorem}
\[
\mathbf{P}[Y=n] = \binom{n-1}{k-1}p^k(1-p)^{n-k}, \qquad n=k,k+1,k+2,\dots
\]
\end{theorem}
\begin{proof}
$\{Y=n\}$ occurs iff exactly $k-1$ successes occur among the first $n-1$ trials (in any of
$\dbinom{n-1}{k-1}$ patterns, each with probability $p^{k-1}(1-p)^{n-k}$ by independence) \emph{and}
trial $n$ is itself a success (probability $p$). Multiplying (again by independence across trials)
and summing over the $\dbinom{n-1}{k-1}$ disjoint patterns:
\[
\mathbf{P}[Y=n] = \binom{n-1}{k-1}p^{k-1}(1-p)^{n-k}\cdot p = \binom{n-1}{k-1}p^k(1-p)^{n-k}.
\qedhere
\]
\end{proof}

\begin{proposition}[Sum of Geometrics representation and the mean]
$Y = X_1+X_2+\cdots+X_k$ where $X_i$ is the number of trials \emph{after} the $(i-1)$-st success
until (and including) the $i$-th success; the $X_i$ are i.i.d.\ Geometric$(p)$. Hence, by linearity
of expectation,
\[
\E{Y} = \sum_{i=1}^k \E{X_i} = \frac{k}{p}.
\]
\end{proposition}
\begin{proof}
That the $X_i$ are i.i.d.\ Geometric$(p)$ follows from the independence of the underlying Bernoulli
trials: the process ``restarts'' after each success, with the same success probability $p$ at
every trial, regardless of the past. Linearity of expectation (proved in Section~6 for two
variables, and extending inductively to any finite number of random variables, discrete or not,
independent or not) then gives $\E{Y}=k\E{X_1}=k/p$.
\end{proof}

\begin{remark}
The variance formula $\Var{Y} = k(1-p)/p^2$ also holds, and follows from $\Var{Y}=\sum_i\Var{X_i}$
\emph{because the $X_i$ are independent}. We do not prove the general fact ``variance of a sum of
\emph{independent} random variables equals the sum of the variances'' here, since it requires the
notion of covariance, which lies beyond the scope of the source notes; it is flagged here only as
a forward reference.
\end{remark}

\subsubsection{Hypergeometric}

Urn with $w$ white, $b$ black balls; draw $n\le w+b$ without replacement; $X=$ number of white
balls drawn.

\begin{theorem}
\[
\mathbf{P}[X=k] = \frac{\dbinom{w}{k}\dbinom{b}{n-k}}{\dbinom{w+b}{n}}, \qquad
\max(0,n-b)\le k \le \min(n,w).
\]
\end{theorem}
\begin{proof}
The sample space (unordered) consists of all $\dbinom{w+b}{n}$ equally likely $n$-subsets of the
$w+b$ balls. The favourable outcomes -- those with exactly $k$ white balls -- are formed by
choosing $k$ of the $w$ white balls ($\dbinom{w}{k}$ ways) \emph{and}, independently, $n-k$ of the
$b$ black balls ($\dbinom{b}{n-k}$ ways); by the basic principle of counting, there are
$\dbinom{w}{k}\dbinom{b}{n-k}$ favourable outcomes. Dividing by the total gives the naive
probability.
\end{proof}

\begin{theorem}
$\E{X} = \dfrac{nw}{w+b}$.
\end{theorem}
\begin{proof}
Label the white balls $1,\dots,w$ and let $Y_i := \mathbf{1}_{\{\text{white ball } i \text{ is
among the } n \text{ drawn}\}}$, so $X = \sum_{i=1}^w Y_i$.

By the same reasoning as Example~\ref{ex:urn}(c) (symmetry among the $\dbinom{w+b}{n}$ equally
likely $n$-subsets): the number of $n$-subsets containing a \emph{particular} ball is
$\dbinom{w+b-1}{n-1}$ (choose the remaining $n-1$ balls from the other $w+b-1$), so
\[
\E{Y_i} = \mathbf{P}[Y_i=1] = \frac{\dbinom{w+b-1}{n-1}}{\dbinom{w+b}{n}} = \frac{n}{w+b}
\]
(the same algebraic simplification as before), for \emph{every} $i=1,\dots,w$ -- notably
\emph{without} needing the $Y_i$ to be independent. By linearity of expectation,
\[
\E{X} = \sum_{i=1}^w \E{Y_i} = w\cdot\frac{n}{w+b} = \frac{nw}{w+b}. \qedhere
\]
\end{proof}

\begin{remark}
$\Var{X} = n\cdot\dfrac{w}{w+b}\cdot\dfrac{b}{w+b}\cdot\dfrac{w+b-n}{w+b-1}$ is the standard result,
but its derivation requires the covariance $\mathrm{Cov}(Y_i,Y_j)$ of the (dependent!) indicators
$Y_i,Y_j$, which is beyond the present scope; we record it here only as a forward reference. Note
the extra factor $\frac{w+b-n}{w+b-1}$ (the ``finite population correction'') relative to the
Binomial variance $n\cdot\frac{w}{w+b}\cdot\frac{b}{w+b}$ one would get from sampling \emph{with}
replacement -- reflecting the negative dependence induced by sampling \emph{without} replacement.
\end{remark}

\subsubsection{Poisson$(\lambda)$}

\begin{theorem}[Poisson as a limit of Binomial]
\label{thm:poisson-limit}
Fix $c>0$. As $n\to\infty$ with $p=c/n$,
\[
\binom{n}{k}\left(\frac{c}{n}\right)^{\!k}\left(1-\frac{c}{n}\right)^{\!n-k}
\longrightarrow \frac{e^{-c}c^k}{k!}, \qquad k=0,1,2,\dots
\]
\end{theorem}
\begin{proof}
Write
\[
\binom{n}{k}\left(\frac{c}{n}\right)^{\!k} = \frac{n(n-1)\cdots(n-k+1)}{k!}\cdot\frac{c^k}{n^k}
= \frac{c^k}{k!}\prod_{j=0}^{k-1}\left(1-\frac{j}{n}\right).
\]
As $n\to\infty$ (with $k$ fixed), each factor $1-j/n \to 1$, so the product $\to 1$, giving
$\binom{n}{k}(c/n)^k \to c^k/k!$.

For the remaining factor, $\left(1-\dfrac{c}{n}\right)^{n-k} =
\left(1-\dfrac{c}{n}\right)^{\!n}\left(1-\dfrac{c}{n}\right)^{\!-k}$. As $n\to\infty$,
$\left(1-\dfrac{c}{n}\right)^{n}\to e^{-c}$ (the standard limit definition of the exponential
function), while $\left(1-\dfrac{c}{n}\right)^{-k}\to 1$ (fixed, finite power, base $\to 1$).

Multiplying the two limits (product of convergent sequences converges to the product of the
limits) gives $\binom{n}{k}(c/n)^k(1-c/n)^{n-k} \to \dfrac{c^k}{k!}\cdot e^{-c}$.
\end{proof}

\begin{definition}
$X\sim\text{Poisson}(\lambda)$ if $\mathbf{P}[X=k] = e^{-\lambda}\lambda^k/k!$, $k\in\Nz$.
\end{definition}

\begin{theorem}
$\E{X}=\lambda$ and $\Var{X}=\lambda$.
\end{theorem}
\begin{proof}
\[
\E{X} = \sum_{k=0}^\infty k\,\frac{e^{-\lambda}\lambda^k}{k!}
= e^{-\lambda}\lambda\sum_{k=1}^\infty \frac{\lambda^{k-1}}{(k-1)!}
= e^{-\lambda}\lambda\, e^{\lambda} = \lambda,
\]
using $\sum_{m=0}^\infty \lambda^m/m! = e^\lambda$ (the exponential series) with $m=k-1$. Similarly,
\[
\E{X(X-1)} = \sum_{k=0}^\infty k(k-1)\frac{e^{-\lambda}\lambda^k}{k!}
= e^{-\lambda}\lambda^2 \sum_{k=2}^\infty \frac{\lambda^{k-2}}{(k-2)!} = e^{-\lambda}\lambda^2 e^\lambda
= \lambda^2,
\]
so $\E{X^2}=\E{X(X-1)}+\E{X}=\lambda^2+\lambda$, and $\Var{X}=\E{X^2}-(\E{X})^2=\lambda^2+\lambda-\lambda^2=\lambda$.
\end{proof}

\subsubsection{Discrete Uniform on $\{1,\dots,n\}$}

\begin{theorem}
If $\mathbf{P}[X=i]=1/n$ for $i=1,\dots,n$, then $\E{X}=\dfrac{n+1}{2}$ and
$\Var{X}=\dfrac{n^2-1}{12}$.
\end{theorem}
\begin{proof}
Using $\sum_{i=1}^n i = \dfrac{n(n+1)}{2}$:
\[
\E{X} = \frac1n\sum_{i=1}^n i = \frac1n\cdot\frac{n(n+1)}{2} = \frac{n+1}{2}.
\]
Using $\sum_{i=1}^n i^2 = \dfrac{n(n+1)(2n+1)}{6}$:
\[
\E{X^2} = \frac1n\sum_{i=1}^n i^2 = \frac{(n+1)(2n+1)}{6}.
\]
Then
\[
\Var{X} = \frac{(n+1)(2n+1)}{6} - \frac{(n+1)^2}{4}
= (n+1)\left[\frac{2(2n+1)-3(n+1)}{12}\right]
= (n+1)\cdot\frac{n-1}{12} = \frac{n^2-1}{12}. \qedhere
\]
\end{proof}

\subsection{Summary comparison table of discrete distributions}

\begin{center}
\renewcommand{\arraystretch}{1.7}
\begin{longtable}{p{2.2cm} p{3.6cm} p{2.6cm} p{2.2cm} p{3.7cm}}
\toprule
\textbf{Distribution} & \textbf{PMF} & \textbf{Mean} & \textbf{Variance} & \textbf{Arises when\dots} \\
\midrule
Bernoulli$(p)$ & $p^x(1-p)^{1-x}$, $x\in\{0,1\}$ & $p$ & $p(1-p)$ & single success/failure trial \\
Binomial$(n,p)$ & $\dbinom{n}{k}p^k(1-p)^{n-k}$, $k=0,\dots,n$ & $np$ & $np(1-p)$ & \# successes in $n$ i.i.d.\ trials (sampling \emph{with} replacement) \\
Geometric$(p)$ & $(1-p)^{k-1}p$, $k=1,2,\dots$ & $1/p$ & $(1-p)/p^2$ & \# trials to \emph{first} success \\
Negative Binomial$(k,p)$ & $\dbinom{n-1}{k-1}p^k(1-p)^{n-k}$, $n=k,k+1,\dots$ & $k/p$ & $k(1-p)/p^2$ & \# trials to $k$-th success \\
Hypergeometric$(w,b,n)$ & $\dbinom{w}{k}\dbinom{b}{n-k}/\dbinom{w+b}{n}$ & $\dfrac{nw}{w+b}$ & $n\dfrac{w}{w+b}\dfrac{b}{w+b}\dfrac{w+b-n}{w+b-1}$ & \# white balls in $n$ draws \emph{without} replacement \\
Poisson$(\lambda)$ & $e^{-\lambda}\lambda^k/k!$, $k=0,1,\dots$ & $\lambda$ & $\lambda$ & limit of Binomial$(n,c/n)$ as $n\to\infty$; rare events \\
Discrete Uniform$\{1,\dots,n\}$ & $1/n$ & $(n+1)/2$ & $(n^2-1)/12$ & a single value chosen ``completely at random'' from a finite set \\
\bottomrule
\end{longtable}
\end{center}

\subsection{Further worked examples}

\begin{example}[Coin tossed until a head or $n$ flips]
Independent flips with $\mathbf{P}[\text{head}]=p$ continue until a head occurs, or until $n$
flips have been made (whichever comes first). Let $X$ be the number of flips made. Find the pmf.

\emph{Solution.} $X$ is discrete, taking values in $\{1,\dots,n\}$. For $k=1,\dots,n-1$:
$\{X=k\}$ means the first $k-1$ flips were tails and the $k$-th was a head, so
$\mathbf{P}[X=k]=(1-p)^{k-1}p$, exactly as in the (untruncated) Geometric distribution. For $k=n$:
$\{X=n\}$ occurs if the first $n-1$ flips are \emph{all} tails (regardless of the $n$-th flip's
outcome, since we stop at $n$ flips regardless), so $\mathbf{P}[X=n]=(1-p)^{n-1}$.

\emph{Check:} $\sum_{k=1}^{n-1}(1-p)^{k-1}p + (1-p)^{n-1}
= p\cdot\dfrac{1-(1-p)^{n-1}}{p} + (1-p)^{n-1} = 1$. \checkmark
\end{example}

\begin{example}[Determining a normalizing constant]
$p(i) = c\lambda^i/i!$ for $i \in \Nz$, with $\lambda>0$ a given constant. Find $c$, then compute
$\mathbf{P}[X=0]$ and $\mathbf{P}[X>2]$.

\emph{Solution.} Normalization requires $\sum_{i=0}^\infty p(i)=1$:
\[
c\sum_{i=0}^\infty \frac{\lambda^i}{i!} = c\,e^{\lambda} = 1 \implies c = e^{-\lambda}.
\]
(This identifies $p(\cdot)$ as exactly the Poisson$(\lambda)$ pmf.) Then
\[
\mathbf{P}[X=0] = p(0) = c = e^{-\lambda},
\]
\[
\mathbf{P}[X>2] = 1-p(0)-p(1)-p(2) = 1 - e^{-\lambda}\left(1+\lambda+\frac{\lambda^2}{2}\right).
\]
\end{example}

\begin{example}[Size-biased sampling: the school-bus problem]
$120$ students ride $3$ buses, with $36, 40, 44$ students respectively. One of the $120$ students
is chosen uniformly at random; let $X$ be the number of students on \emph{that student's} bus.
Find $\E{X}$.

\emph{Solution.} Since the student is chosen uniformly among all $120$ students (not among the $3$
buses), the probability that $X$ equals a given bus's size is proportional to that bus's size:
\[
\mathbf{P}[X=36]=\frac{36}{120}, \quad \mathbf{P}[X=40]=\frac{40}{120}, \quad \mathbf{P}[X=44]=\frac{44}{120}.
\]
So
\[
\E{X} = 36\cdot\frac{36}{120} + 40\cdot\frac{40}{120} + 44\cdot\frac{44}{120}
= \frac{36^2+40^2+44^2}{120} = \frac{1296+1600+1936}{120} = \frac{4832}{120} \approx \boxed{40.27}.
\]
\end{example}

\begin{intuition}
Notice $\E{X}\approx 40.27$ exceeds the simple average bus size $(36+40+44)/3=40$. This is the
\textbf{inspection paradox} / \textbf{size-biased sampling} phenomenon: a randomly chosen
\emph{student} is more likely to be on a large bus than a small one (simply because large buses
carry more students), so the ``bus size as seen by a random student'' systematically overweights
larger buses relative to the plain average over buses.
\end{intuition}

\begin{example}[Betting on high-numbered balls]
Three balls are drawn without replacement from $20$ balls numbered $1$ to $20$. Find the
probability that at least one of the three drawn balls has a number $\ge 17$.

\emph{Solution.} Complement: all three balls have numbers $\le 16$. There are $\dbinom{16}{3}$
ways to choose $3$ balls from the $16$ numbered $1$--$16$, out of $\dbinom{20}{3}$ total ways.
\[
\mathbf{P}(\text{win}) = 1-\frac{\dbinom{16}{3}}{\dbinom{20}{3}} = 1-\frac{560}{1140}
= 1-\frac{28}{57} = \boxed{\frac{29}{57}} \approx 0.5088.
\]
\end{example}

\begin{example}[Random walk / gambler's ruin position]
A particle starts at $0$ on $\Z$ and takes $n$ independent steps, each $+1$ with probability $p$
and $-1$ with probability $1-p$. Let $Y$ be its position after $n$ steps. Find the pmf of $Y$.

\emph{Solution.} Let $X$ = number of $+1$ steps among the $n$; then $X\sim\text{Binomial}(n,p)$,
and the position is $Y = X\cdot(+1) + (n-X)\cdot(-1) = 2X - n$ (a linear/bijective transformation of
$X$ on the relevant range). By LOTUS (or directly, since $Y=2X-n \iff X = (Y+n)/2$),
\[
\mathbf{P}[Y=y] = \mathbf{P}\!\left[X=\frac{y+n}{2}\right]
= \binom{n}{\frac{y+n}{2}}\,p^{\frac{y+n}{2}}(1-p)^{\frac{n-y}{2}},
\]
defined whenever $y+n$ is even and $-n\le y\le n$ (and $0$ otherwise, since $X$ cannot take a
non-integer or out-of-range value).
\end{example}

\newpage
% =========================================================================
\section{Joint Distributions of Random Variables}
% =========================================================================

\subsection{Joint and marginal pmfs}

\begin{definition}
Let $X_1,X_2$ be discrete random variables on the same sample space, with $X_1$ taking values in
$\{a_1,a_2,\dots\}$ and $X_2$ in $\{b_1,b_2,\dots\}$. Their \textbf{joint pmf} is
\[
p(a_i,b_j) := \mathbf{P}[X_1=a_i, X_2=b_j] = \mathbf{P}\big[\{X_1=a_i\}\cap\{X_2=b_j\}\big].
\]
The \textbf{marginal pmf} of $X_1$ is recovered by summing out $X_2$:
\[
p_{X_1}(a_i) = \mathbf{P}[X_1=a_i] = \mathbf{P}\!\left[\{X_1=a_i\}\cap\bigcup_j\{X_2=b_j\}\right]
= \sum_j p(a_i,b_j),
\]
using finite/countable additivity over the disjoint events $\{X_1=a_i,X_2=b_j\}$, $j=1,2,\dots$
(which partition $\{X_1=a_i\}$). Symmetrically $p_{X_2}(b_j)=\sum_i p(a_i,b_j)$.
\end{definition}

\begin{example}[Sum and maximum of two dice]
Roll two distinguishable fair dice; let $X_1$ be their sum, $X_2$ their maximum. For instance,
\[
\mathbf{P}[X_1=7, X_2=5] = \mathbf{P}[\{(2,5),(5,2)\}] = \frac{2}{36} = \frac{1}{18},
\]
since $X_1=7, X_2=5$ forces the (unordered) pair of face values to be $\{2,5\}$, realized by the
two ordered outcomes $(2,5)$ and $(5,2)$ out of the $36$ equally likely outcomes.
\end{example}

\begin{example}[Sum and absolute difference]
Let $U=X_1+X_2$ (sum) and $V=|X_1-X_2|$ (absolute difference) for two fair dice. If $V=0$, then
$X_1=X_2=U/2$, which is only possible if $U$ is even, and corresponds to exactly $1$ of the $36$
outcomes, so $\mathbf{P}[U=u,V=0] = 1/36$ for even $u\in\{2,4,6,8,10,12\}$. If $V=v>0$, then
$\{X_1,X_2\}=\left\{\frac{u+v}{2},\frac{u-v}{2}\right\}$, requiring $u+v$ even and
$1\le \frac{u-v}{2} < \frac{u+v}{2}\le 6$; when these hold there are exactly $2$ ordered outcomes
(swap $X_1,X_2$), so $\mathbf{P}[U=u,V=v]=2/36=1/18$. Otherwise $\mathbf{P}[U=u,V=v]=0$.
\end{example}

\subsection{Independence of random variables}

\begin{definition}
Discrete random variables $X_1,X_2$ (taking values in $\{a_1,a_2,\dots\}$, $\{b_1,b_2,\dots\}$
respectively) are \textbf{independent} if
\[
\mathbf{P}[X_1=a_i,X_2=b_j] = \mathbf{P}[X_1=a_i]\,\mathbf{P}[X_2=b_j] \quad \text{for all } i,j
\]
(equivalently, in terms of cdfs, $\mathbf{P}[X_1\le x_1,X_2\le x_2]=\mathbf{P}[X_1\le
x_1]\mathbf{P}[X_2\le x_2]$ for all $x_1,x_2\in\R$ -- the definition applicable to any, not just
discrete, random variables).
\end{definition}

\begin{example}
Toss a fair coin twice, independently; let $X_i=\mathbf{1}_{\{i\text{-th toss is heads}\}}$. Then
$\mathbf{P}[X_1=1,X_2=1] = \mathbf{P}[HH]=\tfrac14 = \mathbf{P}[X_1=1]\mathbf{P}[X_2=1]$, and
similarly for the other three value-pairs, confirming $X_1,X_2$ are independent -- as one would
expect, since the tosses themselves were assumed independent.
\end{example}

\subsection{Indicator random variables}

\begin{definition}
For an event $A$, the \textbf{indicator random variable} $\mathbf{1}_A$ (also written $X$) is
defined by $\mathbf{1}_A(\omega) = 1$ if $\omega \in A$, and $0$ otherwise.
\end{definition}

Indicator variables are the crucial bookkeeping device that lets us compute expectations of
complicated counting random variables (Binomial, Hypergeometric, etc.) via linearity of
expectation \emph{without} needing to first derive their pmf -- as seen repeatedly above.

\subsection{Linearity of expectation (full proof)}

\begin{theorem}[Linearity of expectation]
\label{thm:linearity}
For any two random variables $X_1,X_2$ on the same sample space (not necessarily independent, and
not necessarily discrete -- though we prove it here in the discrete case),
\[
\E{X_1+X_2} = \E{X_1}+\E{X_2}.
\]
\end{theorem}
\begin{proof}
Let $X_1$ take values in $\{a_1,a_2,\dots\}$, $X_2$ in $\{b_1,b_2,\dots\}$, with joint pmf
$p(\cdot,\cdot)$. By LOTUS applied to the pair $(X_1,X_2)$ and the function $g(a,b)=a+b$:
\begin{equation}
\E{X_1+X_2} = \sum_i\sum_j (a_i+b_j)\,p(a_i,b_j). \tag{$\dagger$}
\end{equation}
By the marginalization identity above, $\sum_j p(a_i,b_j) = \mathbf{P}[X_1=a_i]$ and
$\sum_i p(a_i,b_j)=\mathbf{P}[X_2=b_j]$. Expanding $(\dagger)$ and splitting the double sum:
\[
\E{X_1+X_2} = \sum_i\sum_j a_i\,p(a_i,b_j) + \sum_i\sum_j b_j\,p(a_i,b_j)
= \sum_i a_i\underbrace{\sum_j p(a_i,b_j)}_{=\,\mathbf{P}[X_1=a_i]}
+ \sum_j b_j\underbrace{\sum_i p(a_i,b_j)}_{=\,\mathbf{P}[X_2=b_j]}
\]
(interchanging the order of summation, valid under absolute convergence)
\[
= \sum_i a_i\,\mathbf{P}[X_1=a_i] + \sum_j b_j\,\mathbf{P}[X_2=b_j] = \E{X_1}+\E{X_2}. \qedhere
\]
\end{proof}

\begin{remark}
By induction, $\E{X_1+\cdots+X_n} = \sum_{i=1}^n \E{X_i}$ for any finite collection of random
variables -- \emph{no independence assumption is needed}. This is what licenses the
indicator-variable technique used throughout Section~5 (Binomial and Hypergeometric means).
Independence \emph{is}, however, generally needed for the analogous additivity of
\emph{variances} (flagged as a forward reference at several points above).
\end{remark}

\newpage
% =========================================================================
\section{Moments, Variance, and Their Basic Algebra}
% =========================================================================

\begin{definition}
For $n\in\N$, the \textbf{$n$-th (raw) moment} of $X$ is $\E{X^n}$; for discrete $X$,
$\E{X^n}=\sum_i x_i^n p(x_i)$ (by LOTUS with $g(x)=x^n$). The \textbf{$n$-th central moment} is
$\E{(X-\E{X})^n}$.
\end{definition}

\begin{definition}
The \textbf{variance} of $X$ is $\Var{X} := \E{(X-\E{X})^2}$ (the second central moment); the
\textbf{standard deviation} is $\SD{X}:=\sqrt{\Var{X}}$.
\end{definition}

\begin{proposition}[Computational formula for variance]
\label{prop:var-comp}
$\Var{X} = \E{X^2} - (\E{X})^2$.
\end{proposition}
\begin{proof}
Let $\mu = \E{X}$. Expand $(X-\mu)^2 = X^2 - 2\mu X + \mu^2$. By linearity of expectation
(Theorem~\ref{thm:linearity}, extended to include the constant $\mu^2$ via
$\E{aX+b}=a\E{X}+b$):
\[
\Var{X} = \E{X^2-2\mu X+\mu^2} = \E{X^2} - 2\mu\,\E{X} + \mu^2 = \E{X^2}-2\mu^2+\mu^2
= \E{X^2}-\mu^2. \qedhere
\]
\end{proof}

\begin{proposition}
For constants $a,b$: $\Var{aX+b} = a^2\Var{X}$.
\end{proposition}
\begin{proof}
Let $\mu=\E{X}$; then $\E{aX+b} = a\mu+b$ (linearity of expectation), so
\[
\Var{aX+b} = \E{\big[(aX+b) - (a\mu+b)\big]^2} = \E{\big[a(X-\mu)\big]^2}
= a^2\E{(X-\mu)^2} = a^2\Var{X}. \qedhere
\]
\end{proof}

\subsection{Comparison table: mean and variance formulas}

\begin{center}
\renewcommand{\arraystretch}{1.5}
\begin{longtable}{p{6.5cm} p{5.5cm}}
\toprule
\textbf{Identity} & \textbf{Requires independence?} \\
\midrule
$\E{X_1+X_2} = \E{X_1}+\E{X_2}$ & No \\
$\E{aX+b}=a\E{X}+b$ & No \\
$\Var{X}=\E{X^2}-(\E{X})^2$ & No \\
$\Var{aX+b} = a^2\Var{X}$ & No \\
$\Var{X_1+X_2} = \Var{X_1}+\Var{X_2}$ & \textbf{Yes} (in general there is a cross term
$2\,\mathrm{Cov}(X_1,X_2)$, which vanishes when $X_1,X_2$ are independent) \\
\bottomrule
\end{longtable}
\end{center}

\newpage
% =========================================================================
\section{Continuous Random Variables}
% =========================================================================

\subsection{Definition and the density function}

\begin{definition}
A random variable $X$ is \textbf{continuous} if its cdf $F(x)=\mathbf{P}[X\le x]$ is
\emph{(absolutely) continuous} everywhere -- in particular (since absolute continuity implies
ordinary continuity), $F$ is continuous at every $x \in \R$.
\end{definition}

\begin{proposition}
If $X$ is continuous, $\mathbf{P}[X=x]=0$ for every fixed $x\in\R$.
\end{proposition}
\begin{proof}
By right-continuity of $F$ (Proposition~\ref{prop:cdf-properties}, valid for \emph{any} random
variable), $\lim_{y\to x^+}F(y)=F(x)$. We separately claim $\lim_{y\to x^-}F(y) = F(x^-) :=
\mathbf{P}[X<x]$: for any sequence $y_n \uparrow x$, the events $\{X\le y_n\}$ increase to
$\{X<x\}$ (an outcome satisfies $X(\omega)\le y_n$ for \emph{some} $n$ iff $X(\omega)<x$, since
$y_n\uparrow x$), so by continuity of probability (Theorem~\ref{thm:continuity-prob}),
$F(y_n)=\mathbf{P}[X\le y_n]\to \mathbf{P}[X<x]$.

If $F$ is continuous at $x$, it is in particular left-continuous there, so
$F(x^-) = \lim_{y\to x^-}F(y) = F(x)$. Hence
\[
\mathbf{P}[X=x] = \mathbf{P}[X\le x] - \mathbf{P}[X<x] = F(x)-F(x^-) = 0. \qedhere
\]
\end{proof}

\begin{definition}
Equivalently, $X$ is continuous if there exists a non-negative function $f:\R\to[0,\infty)$, the
\textbf{probability density function (pdf)}, such that
\[
\mathbf{P}[X\in B] = \int_B f(x)\,dx
\]
for every (measurable) set $B\subset\R$.
\end{definition}

\begin{intuition}[Why ``density''?]
Let $\varepsilon>0$ be small, and suppose $f$ is continuous at $x$ (so $f(y)\approx f(x)$ for $y$
near $x$). Then
\[
\mathbf{P}[x-\varepsilon<X<x+\varepsilon] = \int_{x-\varepsilon}^{x+\varepsilon} f(y)\,dy
\approx \int_{x-\varepsilon}^{x+\varepsilon} f(x)\,dy = f(x)\cdot 2\varepsilon.
\]
Reading $\mathbf{P}[x-\varepsilon<X<x+\varepsilon]$ as the probability ``mass'' contained in the
small interval of ``volume'' (length) $2\varepsilon$, we see $f(x) \approx
\dfrac{\text{mass}}{\text{volume}}$ -- exactly the physical notion of a density.
\end{intuition}

\begin{proposition}[Fundamental properties of a pdf]
Any pdf $f$ satisfies: (1) $f(x)\ge 0$ for all $x$; (2) $\displaystyle\int_{-\infty}^\infty
f(x)\,dx = 1$.
\end{proposition}
\begin{proof}
(1) is definitional. (2): take $B=\R$ in the defining property, so
$\int_\R f(x)dx = \mathbf{P}[X\in\R]=\mathbf{P}(S)=1$.
\end{proof}

Since single points have probability zero for continuous $X$, we have, for any $a\le b$,
\[
\mathbf{P}[a<X<b]=\mathbf{P}[a\le X<b]=\mathbf{P}[a<X\le b]=\mathbf{P}[a\le X\le b]
= \int_a^b f(x)\,dx,
\]
and the cdf/pdf are related by
\[
F(x) = \int_{-\infty}^x f(y)\,dy, \qquad F'(x)=f(x) \ \ \text{(by the Fundamental Theorem of Calculus)}.
\]

\begin{remark}
The cdf is defined, and fundamental, for \emph{every} random variable (discrete or continuous or
neither); the pmf applies only to discrete random variables, and the pdf only to continuous ones.
Strategy: to find the distribution of a random variable, first find its cdf; then, if it is
(absolutely) continuous, differentiate to obtain the pdf.
\end{remark}

\subsection{Expectation and variance for continuous random variables}

\begin{definition}
For continuous $X$ with pdf $f$: $\displaystyle \E{X} = \int_{-\infty}^\infty x f(x)\,dx$.
\end{definition}

\begin{proposition}[LOTUS, continuous case]
For $g:\R\to\R$, $\displaystyle \E{g(X)} = \int_{-\infty}^\infty g(x)f(x)\,dx$.
\end{proposition}
\emph{(The proof, for non-negative $g$, mirrors the discrete case with sums replaced by integrals
and a partition of the range of $g$ into level sets; we omit the measure-theoretic details, which
lie outside the present syllabus, and simply record the identity, exactly as in the source notes.)}

\begin{definition}
$\Var{X} := \E{(X-\mu)^2} = \displaystyle\int_{-\infty}^\infty (x-\mu)^2 f(x)\,dx$, where
$\mu=\E{X}$; equivalently (same proof as Proposition~\ref{prop:var-comp}, verbatim, with sums
replaced by integrals),
\[
\Var{X} = \int_{-\infty}^\infty x^2 f(x)\,dx - \left(\int_{-\infty}^\infty x f(x)\,dx\right)^{\!2}.
\]
\end{definition}

All of the algebraic identities of Section~7 -- $\E{aX+b}=a\E{X}+b$, $\Var{aX+b}=a^2\Var{X}$, and
linearity of expectation $\E{\sum_i X_i}=\sum_i \E{X_i}$ (for any, possibly mixed
discrete/continuous, finite collection of random variables) -- continue to hold verbatim in the
continuous setting, since their proofs rest only on the algebra of the expectation operator, not
on discreteness.

\newpage
% =========================================================================
\section{Supplementary Derivations}
% =========================================================================

This section reproduces, with full rigor, two supplementary derivations included in the source
material as auxiliary handouts.

\subsection{Why the Negative Binomial distribution has its name}

Let $X$ be the number of trials required to achieve the $m$-th success, with per-trial success
probability $p$ and $q:=1-p$.

\begin{proposition}
$\mathbf{P}[X=n] = \dbinom{n-1}{m-1}p^m q^{n-m}$ for $n=m,m+1,\dots$ (as already derived in
Section 5.3.3 for general $k$; here we simply relabel $k\to m$ to match the alternate source, and
proceed to explain the terminology).
\end{proposition}

Let $f := n-m$ denote the number of \emph{failures} that occur before the $m$-th success, so
$n = f+m$. Substituting $n=f+m$ into the pmf:
\[
\mathbf{P}[X=f+m] = \binom{f+m-1}{m-1}p^m q^f = p^m\binom{m+f-1}{f}q^f
\]
(using $\dbinom{f+m-1}{m-1}=\dbinom{f+m-1}{f}$, since choosing which $m-1$ of the $f+m-1$ slots are
successes is equivalent to choosing which $f$ are failures).

\begin{proposition}[Generalized binomial series]
For any real $m$ and $|q|<1$,
\[
(1-q)^{-m} = \sum_{f=0}^\infty \binom{m+f-1}{f}q^f.
\]
\end{proposition}
\begin{proof}
This is Newton's generalized binomial theorem, applied with exponent $-m$ and argument $-q$:
\[
(1-q)^{-m} = \sum_{f=0}^\infty \binom{-m}{f}(-q)^f,
\qquad \binom{-m}{f} := \frac{(-m)(-m-1)\cdots(-m-f+1)}{f!}.
\]
Now,
\[
\binom{-m}{f}(-1)^f = \frac{(-1)^f\,m(m+1)\cdots(m+f-1)}{f!}(-1)^f
= \frac{m(m+1)\cdots(m+f-1)}{f!} = \binom{m+f-1}{f}
\]
(the last equality since $m(m+1)\cdots(m+f-1) = \dfrac{(m+f-1)!}{(m-1)!}$ when $m$ is a positive
integer, matching the usual binomial-coefficient formula; the identity extends to general real $m$
by definition of the generalized binomial coefficient). Substituting proves the claim.
\end{proof}

Hence $\mathbf{P}[X=f+m] = p^m\binom{m+f-1}{f}q^f$ is \emph{exactly} $p^m$ times the coefficient of
$q^f$ in the power-series expansion of $(1-q)^{-m}$ -- which is precisely why this family of
distributions is called the \textbf{Negative Binomial} distribution: it is generated by a
binomial-type series with a \emph{negative} exponent.

\subsection{The mean of the Binomial distribution, via differentiation of the Binomial Theorem}

We now give an alternative proof (independent of, and not relying on, indicator variables or
linearity of expectation) of $\E{X}=np$ for $X\sim\text{Binomial}(n,p)$, exactly as presented in
the source handout, with every step justified.

\begin{theorem}
If $\mathbf{P}[X=i]=\dbinom{n}{i}p^i(1-p)^{n-i}$ for $i=0,\dots,n$, then
$\displaystyle\sum_{i=0}^n i\,\mathbf{P}[X=i] = np$.
\end{theorem}
\begin{proof}
Start from the Binomial Theorem, valid as an identity of polynomials in two variables $x,y$:
\[
(x+y)^n = \sum_{r=0}^n \binom{n}{r}x^r y^{n-r}.
\]
Differentiate both sides with respect to $x$ (treating $y$ as a constant; term-by-term
differentiation of a finite sum is always valid):
\[
n(x+y)^{n-1} = \sum_{r=0}^n \binom{n}{r}\,r\,x^{r-1}y^{n-r}
\]
(the $r=0$ term contributes $0$ since it is multiplied by $r=0$). Multiply both sides by $x$:
\[
nx(x+y)^{n-1} = \sum_{r=0}^n r\binom{n}{r}x^r y^{n-r}.
\]
This is an identity valid for \emph{all} real $x,y$; substitute $x=p$, $y=1-p$, so that
$x+y=1$:
\[
np\cdot(1)^{n-1} = \sum_{r=0}^n r\binom{n}{r}p^r(1-p)^{n-r}
\quad\Longrightarrow\quad
np = \sum_{r=0}^n r\,\mathbf{P}[X=r] = \E{X}. \qedhere
\]
\end{proof}

\begin{remark}
This computation is logically independent of, but numerically consistent with, the
indicator-variable proof of $\E{X}=np$ given in Section~5.3.2. It illustrates a second, purely
analytic technique -- ``differentiate the generating identity, then specialize the variables'' --
that recurs throughout probability theory (e.g.\ in the theory of probability and moment
generating functions, a topic beyond the present notes).
\end{remark}

\newpage
% =========================================================================
\section*{Appendix: Master Comparison Table}
\addcontentsline{toc}{section}{Appendix: Master Comparison Table}
% =========================================================================

\begin{center}
\renewcommand{\arraystretch}{1.6}
\begin{longtable}{p{3.4cm} p{3.6cm} p{3.6cm} p{4.0cm}}
\toprule
\textbf{Object} & \textbf{Defined for} & \textbf{Fundamental property} & \textbf{Relation to others} \\
\midrule
pmf $p(x)$ & discrete $X$ & $p(x)\ge 0$, $\sum_x p(x)=1$ & $F(a)=\sum_{x\le a}p(x)$ \\
pdf $f(x)$ & continuous $X$ & $f(x)\ge 0$, $\int_\R f=1$ & $F(x)=\int_{-\infty}^x f$, $F'=f$ \\
cdf $F(x)$ & any $X$ & monotone, right-continuous, $F(\pm\infty)=1,0$ & $\mathbf{P}[a<X\le b]=F(b)-F(a)$ \\
$\mathbf{E}[X]$ (discrete) & $\sum_x x\,p(x)$ & linear: $\mathbf{E}[aX+b]=a\mathbf{E}[X]+b$ & LOTUS: $\mathbf{E}[g(X)]=\sum_x g(x)p(x)$ \\
$\mathbf{E}[X]$ (continuous) & $\int_\R x f(x)dx$ & same linearity & LOTUS: $\mathbf{E}[g(X)]=\int_\R g(x)f(x)dx$ \\
$\operatorname{Var}[X]$ & $\mathbf{E}[(X-\mathbf{E}[X])^2]$ & $=\mathbf{E}[X^2]-(\mathbf{E}[X])^2$ & $\operatorname{Var}[aX+b]=a^2\operatorname{Var}[X]$ \\
\bottomrule
\end{longtable}
\end{center}

\end{document}
